% Packages and bibliography
\documentclass[english,a4paper,12pt,oneside]{scrbook}
\usepackage[utf8]{inputenc}
\usepackage{amsmath, amsthm, amssymb}
\usepackage[english]{babel}
\usepackage{marvosym}
\usepackage{graphics}
\usepackage{hyperref}
\usepackage{graphicx}
\usepackage{float}
\usepackage{mathtools}
\usepackage{calrsfs}
\DeclareMathAlphabet{\pazocal}{OMS}{zplm}{m}{n}
\usepackage[toc,page]{appendix}
\usepackage{xcolor}
\usepackage{cleveref}
\usepackage[backend=bibtex]{biblatex}
\usepackage[nottoc,notlot,notlof]{tocbibind}
\usepackage{tikz-cd}
\addbibresource{Bibliography/bibliography.bib}

% Formatting options
\setlength{\parindent}{0em}
\parskip = 0.35cm \relax
%\renewcommand{\baselinestretch}{1.25}

% Equations and theorems style
\counterwithin*{equation}{section}
\counterwithin*{equation}{subsection}
\renewcommand{\theequation}{%
  \thesection.%
  \ifnum\value{subsection}>0 \arabic{subsection}.\fi
  \arabic{equation}%
}
\newtheoremstyle{break}
  {\topsep}{\topsep}%
  {\upshape}{}%
  {\bfseries}{}%
  {\newline}{}%
\theoremstyle{break}
\newtheorem{thm}[equation]{Theorem}
\newtheorem{cor}[equation]{Corollary}
\newtheorem{lemma}[equation]{Lemma}
\newtheorem{defn}[equation]{Definition}
\newtheorem{prop}[equation]{Proposition}
\newtheorem{ass}[equation]{Assumption}
\newtheorem{pb}[equation]{Problem}
\newenvironment{mproof}[1][\proofname]{%
  \begin{proof}[#1]$ $\par\nobreak\ignorespaces
}{%
  \end{proof}
}
\renewcommand*{\proofname}{Proof}
\theoremstyle{remark}
\newtheorem{obs}[equation]{Observation}
\newtheorem{es}[equation]{Example}

% Useful commands
\newcommand{\mR}{\mathbb{R}}
\newcommand{\cV}{\pazocal{V}}
\newcommand{\ds}{\displaystyle}
\newcommand{\norm}[1]{\left\lVert#1\right\rVert}
\newcommand{\HN}[1]{\norm{#1}_{H}}
\newcommand{\VN}[1]{\norm{#1}_{V}}
\newcommand{\VSN}[1]{\norm{#1}_{V^*}}
\newcommand{\tr}{\text{tr}}
\newcommand{\cc}{\subset\subset}
\newcommand{\emb}{\hookrightarrow}
\newcommand{\ind}[1]{\{\text{ #1 }\}}
\newcommand{\mind}[1]{$#1$}
\newcommand{\cT}{\pazocal{T}}
\newcommand{\id}{\text{Id}}
\newcommand{\te}{\theta}
\newcommand{\Te}{\Theta}
\newcommand{\tred}[1]{\textcolor{red}{#1}}
\newcommand{\dive}{\text{div}}
\newcommand{\weakc}{\rightharpoonup}

\begin{document}

% Titelseite
\pagestyle{empty}       % keine Seitennummer
  \parbox{1.5cm}{\resizebox*{110pt}{!}{\includegraphics{Logos/blau/2015_Logo_TUM_CMYK.pdf}}}\hspace{310pt}%
  \parbox{1.5cm}{\resizebox*{90pt}{!}{\includegraphics{Logos/08_Mathematik/MA_blau/FAK_MA_CMYK.pdf}}}%
\vspace*{1.5cm}
\begin{center}
{\Huge Technical University of Munich}
\\
\vspace*{1.5cm}
{\huge \textsc{Department of Mathematics}}
\\
\vspace*{3cm}
{\Huge \textbf{[Thesis Title]}}
\\
\vspace*{3cm}
{\Large Master's Thesis}\linebreak \\
{\Large von}\linebreak \\
{\Large [First Name Last Name]}\\
\vspace*{3cm}
{\Large 
\begin{tabular}{ll}
Supervisor: & Prof. Dr. [First Name Last Name]\\
Advisor: & [First Name Last Name]\\
Submission Date: & [Day. Month. Year]
\end{tabular}
}
\end{center}
\newpage    % Seitenwechsel

% Seite 2
\vspace*{18cm}
\noindent
I hereby declare that this thesis is my own work and that no other sources have been used except those clearly indicated and referenced.
\\[2cm]
Place, Date\\
original, hand-written signature
\newpage

% vertikaler Leerraum
\vspace*{2.2cm}
\noindent %kein Einzug
{\Huge \textbf{Acknowledgements}} \\
\vspace*{1.6cm} \\
% Seitennummerierung "r�misch
\pagenumbering{roman}
% Kopfzeilen (automatisch erzeugt)
\pagestyle{headings}
[text of acknowledgements]

% Seite 3
\newpage
\section*{German Abstract}
[abstract text]
\section*{English Abstract}
[abstract text]
%Seite 4
\newpage
\tableofcontents  

\pagenumbering{arabic}  % Nummerierung der Seiten in 'arabisch' % neues Kapitel mit Namen "Introduction"

\chapter{Introduction}  \setcounter{page}{1}   % setzt Seitenzaehlung auf 1

\chapter{Continuous problem}

\section{Gateaux differentiability}

\begin{ass}[Geometry assumptions for the shape optimization problem]
\label{ass:geo_sh}
Let $D\subseteq \mR^n$ be a bounded Lipschitz domain, and $\Omega_r \cc D$ also bounded Lipschitz. Call $U_r:=D\setminus \Omega_r$, another bounded Lipschitz domain.
\end{ass}

\begin{defn}[Admissible transformations]

Given $D$, we consider the set $\cT$ as defined in \cref{def:adm}, based on the perturbations $\Te$, also defined in \cref{def:adm}.

We consider $\cT_a:=\cT \cap \{ \tau \in W^{1,\infty}(\mR^n, \mR^n), \norm{\tau - \id}_{W^{1,\infty}(\mR^n, \mR^n)}<C(U_r)\}$, where $C(U_r)$ is given by \cref{thm:ptb_id_lip}.

This ensures that $T(U_r)\cc D$ is also bounded Lipschitz, $T$ the unique Lipschitz continuous representive of any $\tau \in \cT_a$ (see \cref{prop:lipk}).

Thanks to \cref{prop:ptb_id}, a small $\Te$ perturbation of $\tau \in \cT_a$ done in the $W^{1,\infty}(D;\mR^n)$ topology leaves us inside $\cT_a$.

\end{defn}

Note, we need $T$ to have a Lipschitz inverse to conclude $T(U_r)\cc D$: for $x \in D$, we have $0<\delta = \inf_{d \in \partial D}|x-d|\leq \norm{T^{-1}}_{W^{1,\infty}(\mR^n,\mR^n)}\inf_{d \in \partial D}|T(x)-d|$.

We are interested in solving the following problem.

\begin{pb}[Shape optimization problem]
\label{pb:shopt}
Suppose \cref{ass:geo_sh}, \cref{ass:diri}, \cref{ass:neu} hold. Using the notation $T = U(\tau)$ for "unique continuous representative" of $\tau$, we want to solve:

$$\inf_{\tau \in \cT_a}\frac{1}{2}\norm{v_\tau-w_\tau}_{L^2(I,H_T)}^2=:J(\tau)$$

Here $v^\tau = v_0^\tau + \bar{u}$, and $v_0^\tau, w_\tau$ are given in \cref{pb:joint_mov} and $\bar{u}$ is given in \cref{pb:diri_ext}. The notation for the spaces also comes from  \cref{pb:joint_mov}, where $\cdot_T$ means a space defined on the moving domain, $H=L^2, W=H^1_0, V = H^1_c$ (and the last space is defined in \cref{subs:inh_diri}).

\end{pb}

\begin{obs}[Well posedness of $J$]

We know from \cref{prop:diri_wp} that $v_0^T + \bar{u}$ doesn't depend on the particular choice of $\bar{u}$, therefore, for different $\tau$'s yielding the same domain, $J(\tau)$ doesn't change.

\end{obs}

Moreover, \cref{ass:geo_sh} ensures in particular \cref{ass:pull}, which enables the applicability of \cref{thm:eq_pde}.

Now, let $U:=T(U_r)$, for $T \in \cT_a$ (we write $T \in \cT_a$ for $\tau \in \cT_a$, $U(\tau)=T$) and let $\delta \te \in \Te$. 

We are interested in finding $J'(\tau) \in \Te^*$ such that, for all $\delta \te_k \rightarrow 0$ in $\Te$, we have:


$$\lim_{k}\frac{|J(\tau+\delta \te_k)-J(\tau)-J'(\tau)(\delta \te_k)|}{\norm{\delta \te_k}_{\Te}}$$

For simplificity we have put $\norm{\te}_\Te = \norm{\te}_{W^{1,\infty}(\mR^n;\mR^n)}=\norm{\te}_{W^{1,\infty}(D;\mR^n)}$.

Note, $\tau+\delta \te_k \in \cT_a$ for large enough $k$. In fact, $\tau+\delta \te_k \in \cT$ for large $k$ as in \cref{prop:ptb_id}, and the condition on $C(U_r)$ is satisfied too, because $\te :=\tau-\id$ was already in the $W^{1,\infty}$ open ball of radius $C(U_r)$ centered at the origin, and so will be $\te + \delta_k\te$, always for large $k$.

Now, $\tau+\delta \te_k  = (\id+\delta\te_k \circ \tau^{-1})\circ \tau$, and $\id+\delta\te_k \circ \tau^{-1}$ is in $\cT_a$ (it is in $\cT$ by \cref{prop:ptb_id} and the reasoning above shows it is also in $\cT_a$).

We are then equivalently interested in:

$$\lim_{k}\frac{|J((\id+\delta\te_k \circ \tau^{-1})\circ \tau)-J(\tau)-J'(\tau)(\delta \te_k)|}{\norm{\delta \te_k}_{\Te}}$$

This amounts to setting the reference domain to $T(U_r)$ instead of $U_r$ and perturbing the former, at least for the sake of computing derivatives.

We now introduce a Lagrangian.

For $k\geq K(T)$ we define $T_k:=\id+\delta\te_k \circ \tau^{-1} \in \cT_a$, as seen above.

With the help and the notation of \cref{thm:eq_pde} we can set:

\begin{align*}
L_T(k,w,v_0,q,p) = \\
\frac{1}{2}\int_I \int_{T(U_r)}|v_0+\bar{u}_T\circ T_k - w|^2|\det(DT_k)|+\\
\int_I ( w_t , q |\det(DT_k)|)_{H_T}+ (A_{T_k}\nabla w, \nabla q)_{H_T} -\int_I(g,\tr_{U} q)_{L^2(\Gamma_f)} +\\ \int_I (v_{0t},p |\det(DT_k)|)_{H_T} + (A_{T_k} \nabla v_0, \nabla p)_{H_T}+\\\int_I((\bar{u}\circ T_k)',p|\det(DT_k)|)_{H_T}+(A_{T_k} \nabla (\bar{u} \circ T_k), \nabla p)_{H_T}
\end{align*}

for $w \in Q_0(I, V_T), v_0 \in Q_0(I,W_T), q \in Q^0(I, V_T), p \in Q^0(I, W_T)$. Remember,  $A_T:=  (DT)^{-1}(DT)^{-t}|\det(DT)|$.

Note that to be precise, $\bar{u}$ is an extension (any extension in fact, satisfying the conditions of \cref{pb:diri_ext}) of the Dirichlet datum $f$, on the moving domain $T_k(T(U_r))$. Because of this, let's fix $\bar{u}_T$ with this property on $T(U_r)$. We show that $\bar{u}_T\circ T_k^{-1}$ satisfies the conditions stated in \cref{pb:diri_ext}. 

In particular:

\begin{itemize}
	\item composition with $T$ preserves the smoothness of the extension, as seen in \cref{lemma:bochner_Hk_map}, given that $\circ T_k^{-1}$ is a linear bounded operator between $V_T$ and $V_{T_k \circ T}$ (see \ref{thm:change})
	\item the initial value is preserved, as seen in the proof of \cref{prop:change_boch}
	\item the trace on $\Sigma_f$ is preserved, because the trace on $\Gamma_f=\partial D$ is preserved, see \cref{thm:change}
\end{itemize}

Therefore we can state the following definition.

\begin{defn}[Lagrangian]

For a fixed $T \in \cT_a$ and  $k\geq K(T)$, for $T_k:=\id+\delta\te_k \circ \tau^{-1} \in \cT_a$, we define:

\begin{align*}
L_T(k,w,v_0,q,p) = \\
\frac{1}{2}\int_I \int_{T(U_r)}|v_0+\bar{u}_T - w|^2|\det(DT_k)|+\\
\int_I ( w_t , q |\det(DT_k)|)_{H_T}+ (A_{T_k}\nabla w, \nabla q)_{H_T} -\int_I(g,\tr_{U} q)_{L^2(\Gamma_f)} +\\ \int_I (v_{0t},p |\det(DT_k)|)_{H_T} + (A_{T_k} \nabla v_0, \nabla p)_{H_T}+\\\int_I(\bar{u}_T',p|\det(DT_k)|)_{H_T}+(A_{T_k} \nabla \bar{u}_T , \nabla p)_{H_T}
\end{align*}

$L_T$ is defined as a map $\{k\geq K(T)\}\times Q_0(I, V_T)\times Q_0(I,W_T)\times Q^0(I, V_T)\times Q^0(I, W_T)\rightarrow \mR$.

We call $u = (w,v_0)$, $\pi = (q,p)$, $G(k,u,\pi) = L_T(k,w,v_0,q,p)$ to ease the notation.

We also call $b(k, u) = \frac{1}{2}\int_I \int_{T(U_r)}|v_0+\bar{u}_T - w|^2|\det(DT_k)|$ and $a(k, u,\pi) = G(k,u,\pi)-b(k, u)$, $E = Q_0(I, V_T)\times Q_0(I,W_T)$, $F=Q^0(I, V_T)\times Q^0(I, W_T)$.

\end{defn}

\begin{prop}[Properties of the Lagrangian]
\label{prop:lagr}

$L_T$ satisfies the following properties:

\begin{enumerate}
	\item $\psi \mapsto a(k, \phi,\psi)$ is linear, no matter what $\phi,k$
	\item $G$ is Frechet differentiable with respect to $\psi$ at $(k,\phi,0)$ for all $k, \phi$
	\item $d_\psi G(k,\phi,0)[\delta \psi]=0$ for all $\delta \psi \in F$ admits a unique solution $\phi = u^k$
	\item $[0,1]\ni s \mapsto G(k, su^k + (1-s)u^0,\psi)$ is $AC[0,1]$, no matter what $k, \psi$
	\item $G$ is Frechet differentiable with respect to $\phi$ at $(k,\psi,\phi)$ for all $k, \psi, \phi$
	\item $[0,1]\ni s \mapsto d_\phi G(k, su^k + (1-s)u^0,\psi)[\delta \phi]$ is $L^1(0,1)$, no matter what $k, \psi, \delta \phi$
	\item there exists a unique solution $\psi = \pi^k$ to $\ds \int_0^1 d_\phi G(k, su^k + (1-s)u^0,\psi)[\delta \phi]ds =0$ for all $\delta \psi$ 
\end{enumerate}

In particular $\pi^k = (Q^k \circ \tau^k,P^k \circ \tau^k)$, where we introduced the averaged adjoint problems on the moving domain:

\begin{pb}[Averaged adjoint equations]
\label{pb:avg_adj_pb}
\begin{align*}
-Q^k_t-\Delta Q^k =\frac{v_0^k-w^k+v_0^0-w^0}{2}\circ T_k^{-1}+\bar{u}_T\circ T_k^{-1} \\
Q^k(T)=0\\
\partial_\nu Q^k = 0 \text{ on } \Sigma_f\\
Q^k = 0 \text{ on } \Sigma_m^{T_k\circ T}
\end{align*}

and

\begin{align*}
-P^k_t-\Delta P^k =-\frac{v_0^k-w^k+v_0^0-w^0}{2}\circ T_k^{-1}-\bar{u}_T\circ T_k^{-1} \\
P^k(T)=0\\
P^k = 0 \text{ on } \Sigma_m^{T_k\circ T} \sqcup \Sigma_f
\end{align*}
\end{pb}

\end{prop}

\begin{mproof}

The first point is immediate.

\underline{Proof of 2}

All the pieces are linear in $\psi$. We only check the boundedness of the various differentials. For simplicity, call $|\det(DT_k)|=d$, and note that   $\norm{qd}_{H_T}\leq C(d)\norm{q}_{H_T}$.

And now, for instance:
%\int_I \langle w_t , q |\det(DT_\e)|\rangle_{V^*_T,V_T}+ (A_{T_\e}\nabla w, \nabla q)_{H_T} -\int_I(g,\tr_{U} q)_{L^2(\Gamma_f)}
\begin{align*}
\int_I ( w_t , \delta q |\det(DT_k)|)_{H_T} = \int_I ( w_t , \delta q d )_{H_T}\leq\\ C(d) \int_I \norm{w_t}_{H_T}\norm{\delta q}_{H_T}\leq C(d) \norm{w_t}_{L^2(I,H_T)}\norm{\delta q}_{L^2(I,H_T)}\leq\\C(d) \norm{w}_{Q(I,V_T)}\norm{\delta q}_{Q(I,V_T)}\leq
C(d) \norm{w}_{Q(I,V_T)}(\norm{\delta q}_{Q(I,V_T)}+\norm{\delta p}_{Q(I,V_T)}) =\\ C(d) \norm{w}_{Q(I,V_T)}\norm{\delta\psi}_F
\end{align*}

Or also:

\begin{align*}
\int_I(g,\tr_{U} \delta q)_{L^2(\Gamma_f)}\leq \int_I \norm{g}_{L^2(\Gamma_f)}\norm{\delta q}_{V_T}\leq \norm{g}_{H^1(I,L^2(\Gamma_f))}\norm{\delta\psi}_F
\end{align*}

and:

\begin{align*}
\int_I (A_{T_k} \nabla \bar{u}_T, \nabla\delta p)_{H_T}\leq C\norm{A_{T_k}}_{L^\infty(D;\mR^{n\times n})}\int_I \norm{\nabla \bar{u}_T}_{L^2(I,H_T)}\norm{\nabla\delta p}_{L^2(I,H_T)}\leq\\ C(T) \norm{\delta\psi}_F \norm{\bar{u}}_{H^1(I,V_T)}
\end{align*}

\underline{Proof of 3}

We get back the state equations, thanks to linearity, and by testing separately with $\delta \psi =(\delta q, 0)$ and $\delta \psi = (0,\delta p)$, so that a unique solution exists by \cref{thm:eq_pde}.

\underline{Proof of 4}

Every piece but $b$ is linear or constant in the state $\phi$. We only need to prove that $[0,1]\ni s \mapsto b(k, su^k + (1-s)u^0)$ is $AC[0,1]$. But by the structure of the cost function $J$, transported on $T(U_r)$, we see that the latter is a quadratic polynomial in $s$, hence, absolutely continuous.

\underline{Proof of 5}

For the pieces with the gradients, it follows as above, by in case employing the simmetry of $A_{T_k}$.

Now, for instance the linear form $\delta v_0 \mapsto \int_I (\delta v_{0t},p |\det(DT_k)|)_{H_T}$ is also bounded by $C(d) \norm{\delta v_0}_{Q(I,W_T)}\norm{\delta q}_{Q(I,V_T)}$ just like before.

What remains to check is the Frechet differentiability of $b$.

To do so, perturb $\phi$ by $\delta \phi$ and expanding the square:

\begin{align*}
\frac{1}{2}\int_I \int_{T(U_r)}|v_0+\delta v_0+\bar{u}_T - w-\delta w|^2|\det(DT_k)| = \\\frac{1}{2}\int_I \int_{T(U_r)}|v_0+\bar{u}_T - w|^2|\det(DT_k)|+\\\frac{1}{2}\int_I \int_{T(U_r)}|\delta v_0-\delta w|^2|\det(DT_k)|+\\\int_I \int_{T(U_r)}(v_0+\bar{u}_T - w)(\delta v_0-\delta w)|\det(DT_k)|
\end{align*} 

Now, $\ds \int_I \int_{T(U_r)}|\delta v_0-\delta w|^2|\det(DT_k)|\leq C(T_k)\norm{\delta v_0-\delta w}_{L^2(I,H_T)}^2\leq C(T_k)\norm{\phi}_E^2$, so that this term is of higher term.

And $\int_I \int_{T(U_r)}(v_0+\bar{u}_T - w)(\delta v_0-\delta w)|\det(DT_k)|$ is linear and bounded by reasonings similar to the former ones.

\underline{Proof of 6}

By the last point:

\begin{align*}
d_\phi G(k, \phi ,\psi)[\delta \phi] =\\
\int_I ((v_0+\bar{u}_T - w)|\det(DT_k)|,\delta v_0-\delta w)_{H_T}+\\
\int_I (\delta w_t , q |\det(DT_k)|)_{H_T}+ (A_{T_k}\nabla \delta w, \nabla q)_{H_T}+\\
\int_I (\delta v_{0t},p |\det(DT_k)|)_{H_T} + (A_{T_k} \nabla \delta v_0, \nabla p)_{H_T}
\end{align*}

so that:

\begin{align*}
d_\phi G(k, su^k + (1-s)u^0,\psi)[\delta \phi] = \\
\int_I ((s(v_0^k+\bar{u}_T - w^k)+(1-s)(v_0^0+\bar{u}_T - w^0))|\det(DT_k)|,\delta v_0-\delta w)_{H_T}+\\
\int_I ( \delta w_t , q |\det(DT_k)|)_{H_T}+ (A_{T_k}\nabla \delta w, \nabla q)_{H_T}+\\
\int_I ( \delta v_{0t},p |\det(DT_k)|)_{H_T} + (A_{T_k} \nabla \delta v_0, \nabla p)_{H_T}
\end{align*}

which is a degree $1$ polynomial in $s$, hence, $L^1(0,1)$.

\underline{Proof of 7}

Rewriting the formula above and integrating in $s$, we come to:

\begin{align*}
\int_0^1d_\phi G(k, su^k + (1-s)u^0,\psi)[\delta \phi]ds = \\
\int_I (((v_0^k+\bar{u}_T - w^k)+(v_0^0+\bar{u}_T - w^0))/2|\det(DT_k)|,\delta v_0-\delta w)_{H_T}+\\
\int_I ( \delta w_t , q |\det(DT_k)|)_{H_T}+ (A_{T_k}\nabla \delta w, \nabla q)_{H_T}+\\
\int_I ( \delta v_{0t},p |\det(DT_k)|)_{H_T} + (A_{T_k} \nabla \delta v_0, \nabla p)_{H_T}
\end{align*}

As in \cref{prop:change_boch}, $ \delta w_t  = (\delta w\circ T_k^{-1})_t\circ T_k$, where $\delta w\circ T_k^{-1} \in Q_0(I,V_{T_k \circ T})$ by \cref{prop:change_boch} (that can be applied thanks to the smallness of $T_k$).

Applying a change of variables we are left with:

\begin{align*}
\int_0^1 d_\phi G(k, su^k + (1-s)u^0,\psi)[\delta \phi]ds = \\
\int_I \left (\frac{v_0^k-w^k}{2}\circ T_k^{-1}+ \frac{v_0^0-w^0}{2}\circ T_k^{-1}+\bar{u}_T\circ T_k^{-1} ,\delta v_0\circ T_k^{-1}-\delta w\circ T_k^{-1}\right)_{H_{T_k \circ T}}+\\
\int_I ((\delta w\circ T_k^{-1})_t , q\circ T_k^{-1} )_{H_{T_k \circ T}}+ (\nabla (\delta w\circ T_k^{-1}), \nabla( q\circ T_k^{-1}))_{H_{T_k \circ T}}+\\
\int_I ( (\delta v_{0}\circ T_k^{-1})_t,p \circ T_k^{-1})_{H_{T_k \circ T}} + ( \nabla \delta (v_0\circ T_k^{-1}), \nabla (p\circ T_k^{-1}))_{H_{T_k \circ T}}
\end{align*}

Here, as we saw in \cref{prop:change_boch}, we have $\delta w\circ T_k^{-1}, w\circ T_k^{-1} \in Q_0(I, V_{T_k \circ T})$, $ \delta v_{0}\circ T_k^{-1}, v_{0}\circ T_k^{-1} \in Q_0(I,W_{T_k \circ T})$, $q\circ T_k^{-1}\in Q^0(I, V_{T_k \circ T})$ and $p\circ T_k^{-1}\in Q^0(I, W_{T_k \circ T})$.

Because $\circ T_k^{-1}$ is a bijection of $Q_0(I,W_{T_k \circ T})$ and $Q_0(I,W_{T_k})$ as we saw in \ref{prop:change_boch} (and analogously of $V$), we have that  $\int_0^1 d_\phi G(k, su^k + (1-s)u^0,\psi)[\delta \phi]ds=0$ for all $\delta \phi \in E$ if and only if:

\begin{align*}
\int_I \left (\frac{v_0^k+w^k}{2}\circ T_k^{-1}- \frac{v_0^0+w^0}{2}\circ T_k^{-1}+\bar{u}_T\circ T_k^{-1} ,\delta V_0-\delta W\right)_{H_{T_k \circ T}}+\\
\int_I ( \delta W_t , q\circ T_k^{-1})_{H_{T_k \circ T}}+ (\nabla\delta W, \nabla( q\circ T_k^{-1}))_{H_{T_k \circ T}}+\\
\int_I ( \delta V_{0t},p \circ T_k^{-1})_{H_{T_k \circ T}} + ( \nabla \delta V_0, \nabla (p\circ T_k^{-1}))_{H_{T_k \circ T}} = 0
\end{align*}

for all $\delta W, \in Q_0(I, V_{T_k \circ T})$, $ \delta V_{0} \in Q_0(I,W_{T_k \circ T})$.

We wish to find a (unique) solution $(q^k, p^k) \in Q^0(I, V_T)\times Q^0(I, W_T)$ of this problem. We can equivalently (by \cref{prop:change_boch}) find $(Q^k, P^k) \in Q^0(I, V_{T_k \circ T})\times Q^0(I, W_{T_k \circ T})$ satisfying:

\begin{align*}
\int_I \left (\frac{v_0^k-w^k}{2}\circ T_k^{-1}+ \frac{v_0^0-w^0}{2}\circ T_k^{-1}+\bar{u}_T\circ T_k^{-1} ,\delta V_0-\delta W\right)_{H_{T_k \circ T}}+\\
\int_I (\delta W_t ,Q^k )_{H_{T_k \circ T}}+ (\nabla \delta W, \nabla Q^k)_{H_{T_k \circ T}}+\\
\int_I( \delta V_{0t},P^k)_{H_{T_k \circ T}} + ( \nabla \delta V_0, \nabla P^k)_{H_{T_k \circ T}} = 0
\end{align*}

for all $\delta W, \in Q_0(I, V_{T_k \circ T})$, $ \delta V_{0} \in Q_0(I,W_{T_k \circ T})$.

By testing first with $\delta W=0$ and then with $\delta V_0=0$ we can equivalently look for:

\begin{align*}
(Q^k, P^k) \in Q^0(I, V_{T_k \circ T})\times Q^0(I, W_{T_k \circ T}) \text{ with }\\
\int_I (\delta W_t ,Q^k)_{H_{T_k \circ T}}+ (\nabla \delta W, \nabla Q^k)_{H_{T_k \circ T}} =\\ \int_I \left (\frac{v_0^k+w^k-v_0^0-w^0}{2}\circ T_k^{-1}+\bar{u}_T\circ T_k^{-1} ,\delta W\right)_{H_{T_k \circ T}}\\
\int_I ( \delta V_{0t},P^k)_{H_{T_k \circ T}}+ ( \nabla \delta V_0, \nabla P^k)_{H_{T_k \circ T}} =\\- \int_I \left (\frac{v_0^k-w^k+v_0^0-w^0}{2}\circ T_k^{-1}+\bar{u}_T\circ T_k^{-1} ,\delta V_0\right)_{H_{T_k \circ T}}
\end{align*}

An application of integration by parts in time (see \cref{prop:Q}) yields the problem:

\begin{align*}
(Q^k, P^k) \in W^0(I, V_{T_k \circ T})\times W^0(I, W_{T_k \circ T}) \text{ with }\\
-\int_I ( Q^k_t, \delta W )_{H_{T_k \circ T}}+ (\nabla \delta W, \nabla Q^k)_{H_{T_k \circ T}} =\\ \int_I \left (\frac{v_0^k-w^k+v_0^0-w^0}{2}\circ T_k^{-1}+\bar{u}_T\circ T_k^{-1} ,\delta W\right)_{H_{T_k \circ T}}\\
-\int_I( P^k_t, \delta V_0)_{H_{T_k \circ T}} + ( \nabla \delta V_0, \nabla P^k)_{H_{T_k \circ T}} =\\- \int_I \left (\frac{v_0^k-w^k+v_0^0-w^0}{2}\circ T_k^{-1}+\bar{u}_T\circ T_k^{-1} ,\delta V_0\right)_{H_{T_k \circ T}}
\end{align*}

But this is the weak formulation (cfr. \cref{thm:eq_pde}, \cref{pb:diri_ext}, \cref{pb:neu}) of the problems:

\begin{align*}
-Q^k_t-\Delta Q^k =\frac{v_0^k-w^k+v_0^0-w^0}{2}\circ T_k^{-1}+\bar{u}_T\circ T_k^{-1} \\
Q^k(T)=0\\
\partial_\nu Q^k = 0 \text{ on } \Sigma_f\\
Q^k = 0 \text{ on } \Sigma_m
\end{align*}

and

\begin{align*}
-P^k_t-\Delta P^k =-\frac{v_0^k-w^k+v_0^0-w^0}{2}\circ T_k^{-1}-\bar{u}_T\circ T_k^{-1} \\
P^k(T)=0\\
P^k = 0 \text{ on } \Sigma_m \sqcup \Sigma_f
\end{align*}

Applying the time reversal $t\mapsto T -t$ (where $I = [0,T]$), these are a couple of standard heat equations for which we have available existence, uniqueness and stability results (see \cref{chap:parab_eq}, and \cref{prop:eq_form}).

By calling then $\pi^k = (Q^k \circ T^k,P^k \circ T^k)$ we conclude the proof.

\end{mproof}

We now turn to the verification of Gateaux differentiability of $J$, applying the techniques proposed in \cite{avg_adj}.

\begin{prop}[Averaged adjoint method for Gateaux derivatives]
\label{prop:adv_adj}

If $J'(\tau) \in \Te^*$ satisfies:

$$\lim_{k}\frac{G(k,u^0,\pi^k)-G(0,u^0,\pi^k)}{t_k}=J'(\tau)[\delta \te]$$

where $\delta\te_k = t_k\delta \te$ for $t_k\rightarrow 0$, then $J'(\tau)$ is the Gateaux derivative of $J$ at $\tau$.

\end{prop}

\begin{mproof}

We have $G(k,u^k,\pi^k)-G(k,u^0,\pi^k)  = \int_0^1 d_\phi G(k, su^k + (1-s)u^0,\pi^k)[u^k-u^0]ds = 0$ because $u^k-u^0 \in E$, and by absolute contintinuity and integrability of derivative as seen in \cref{prop:lagr}.

Moreover, calling $g_k = G(k,u^k,0)-G(0,u^0,0)$, we have:

\begin{itemize}
	\item $g_0 = 0$
	\item $g_k = J((\id+\delta\te_k \circ \tau^{-1})\circ \tau)-J(\tau)$, thanks again to a change of variables
	\item $g_k = G(k,u^k,\pi^k)-G(0,u^0,\pi^k)$ thanks to $\pi^k \in F$ and the state equations (note, this is possible because $k,u^k$ appear, and $0,u^0$ appear, so that the indices don't mix)
\end{itemize}

And now, $J(\tau+\delta\te_k)-J(\tau) = J((\id+\delta\te_k \circ \tau^{-1})\circ \tau)-J(\tau) = g_k = G(k,u^k,\pi^k)-G(0,u^0,\pi^k) = G(k,u^k,\pi^k)-G(k,u^0,\pi^k)+G(k,u^0,\pi^k)-G(0,u^0,\pi^k) = G(k,u^0,\pi^k)-G(0,u^0,\pi^k)$.

\end{mproof}

\begin{prop}[Gateaux differentiability of $J$]
Given $\tau \in \cT_a$, $J$ is Gateaux differentiable at $\tau$ with respect to the $W^{1,\infty}$ topology. The Gateaux differential is:


\begin{align*}
J'(\tau)[\delta \te] =\\ \int_I (w_t^\tau \dive(\delta \te\circ  \tau^{-1}), q^\tau )+ \int_I (A'(\delta\te \circ \tau^{-1})\nabla v^\tau, \nabla p^\tau)+\\
\int_I (v_t^\tau \dive(\delta \te\circ  \tau^{-1}), p^\tau )+ \int_I (A'(\delta\te \circ \tau^{-1})\nabla w^\tau, \nabla q^\tau)+\\
\frac{1}{2}\int_I\int_{T(U_r)}|v^\tau-w^\tau|^2\dive(\delta \te\circ  \tau^{-1})
\end{align*}

where $p^\tau$, $q^\tau$ solve:

\begin{align*}
-q^\tau_t-\Delta q^\tau =v^\tau-w^\tau\\
q^\tau(T)=0\\
\partial_\nu q^\tau = 0 \text{ on } \Sigma_f\\
q^\tau = 0 \text{ on } \Sigma_m^{T_k\circ T}
\end{align*}

and

\begin{align*}
-p^\tau_t-\Delta p^\tau = - v^\tau+ w^\tau \\
p^\tau(T)=0\\
p^\tau = 0 \text{ on } \Sigma_m^{T_k\circ T} \sqcup \Sigma_f
\end{align*}

and where $A'(\delta\te )= - D\delta \te -(D\delta\te)^t + \dive(\delta\te)$.

\end{prop}

\begin{mproof}

\underline{The shape derivative is linear and bounded}

Linearity is immediate. For the boundedness:

\begin{align*}
|J'(\tau)[\delta \te]| \leq\\ \norm{\dive(\delta \te\circ  \tau^{-1})}_{L^\infty(T(U_r))}\left (\int_I( \norm{q_t^\tau}_{H^T}\norm{q^\tau}_{H_T}+ \norm{v_t^\tau}_{H_T}\norm{p^\tau}_{H_T})+\frac{1}{2}\norm{v^\tau-w^\tau}^2_{L^2(I,H_T)}\right )+\\
\left(\sum_{ij} \norm{(A'(\delta\te \circ \tau^{-1})_{ij}}_{L^\infty(T(U_r))}\right )\left (
\int_I \norm{\nabla v^\tau}_{H_T} \norm{\nabla p^\tau}_{H_T}+ \int_I \norm{\nabla w^\tau}_{H_T} \norm{\nabla q^\tau}_{H_T}\right )
\end{align*}

and then, for $C$ independent of $\delta \te$:

\begin{align*}
|J'(\tau)[\delta \te]| \leq\\ C \left ( \norm{\dive(\delta \te\circ  \tau^{-1})}_{L^\infty(T(U_r))}+\left(\sum_{ij} \norm{(A'(\delta\te \circ \tau^{-1})_{ij}}_{L^\infty(T(U_r))}\right )\right )\leq\\
\norm{\delta \te \circ \tau^{-1}}_{W^{1,\infty}(\mR^n;\mR^n)}\leq \\
C\norm{\delta \te }_{W^{1,\infty}(\mR^n;\mR^n)}
\end{align*}

where in the last step we applied point i) of lemme 2.2, \cite{murat}. This shows the boundedness.

\underline{Conclusion}

Assume $p^k \weakc p^0$ in $Q(I,W_T)$ and $q^k \weakc q^0$ in $Q(I,V_T)$.

Now, using that $u^0=(w^\tau, v_0^\tau)$:

\begin{align*}
G(k,u^0,\pi^k)-G(0,u^0,\pi^k) =\\
\frac{1}{2}\int_I \int_{T(U_r)}|v^\tau-w^\tau|^2|\det(DT_k)|+\\
\int_I ( w_t^\tau |\det(DT_k)| , q^k)_{H_T}+ (A_{T_k}\nabla w^\tau, \nabla q^k)_{H_T} -\int_I(g,\tr_{U} q^k)_{L^2(\Gamma_f)} +\\ \int_I (v_t^\tau |\det(DT_k)|,p^k )_{H_T} + (A_{T_k} \nabla v^\tau, \nabla p^k)_{H_T} - \\
\frac{1}{2}\int_I \int_{T(U_r)}|v^\tau-w^\tau|^2-\\
\int_I ( w_t^\tau , q^k )_{H_T}+(\nabla w^\tau, \nabla q^k)_{H_T} +\int_I(g,\tr_{U} q^k)_{L^2(\Gamma_f)} -\\ \int_I (v_t^\tau,p^k )_{H_T} + ( \nabla v^\tau, \nabla p^k)_{H_T} \\
\end{align*}

Grouping some terms and cancelling the boundary integral:

\begin{align*}
G(k,u^0,\pi^k)-G(0,u^0,\pi^k) =\\
\frac{1}{2}\int_I \int_{T(U_r)}|v^\tau-w^\tau|^2(|\det(DT_k)|-1)+\\
\int_I ( w_t^\tau (|\det(DT_k)| -1), q^k)_{H_T}+ ((A_{T_k}-I)\nabla w^\tau, \nabla q^k)_{H_T}+\\
\int_I (v_t^\tau (|\det(DT_k)|-1),p^k )_{H_T} + ((A_{T_k}-I) \nabla v^\tau, \nabla p^k)_{H_T} \\
\end{align*}

Now, the application $\delta \te \mapsto \id +\delta \te \circ \tau^{-1}$ is Frechet differentiable at $\delta \te =0$, as a map of $\Te$ into $\cV^1$, with Frechet derivative $\delta \te \circ \tau^{-1}$, which is linear and bounded by point i) of lemme 2.2, \cite{murat}. Note, we needed here $\tau \in \cT^1$.

Also, the maps $\delta \eta \mapsto |\det(D\eta)|$ and $\eta\mapsto (D\eta)^{-1}(D\eta)^{-t}|\det D\eta|$ are Frechet differentiable at $\id$, from $\cV^1$ into $L^\infty(\mR^n;\mR)$ and $L^\infty(\mR^n;\mR^{n\times n})$, as stated in lemma 4.16, page 80 of \cite{lindemann}. Their Frechet derivatives are $\dive (\beta)$ and $I-D\beta-(D\beta)^t$, respectively.

Therefore, composition with  $\delta \te  \mapsto \id+ \delta \te \circ \tau^{-1}$ yields two Frechet differentiable maps, whose derivatives at $0$, in direction $\delta \te \in \Te$ are exactly:

\begin{itemize}
	\item $\dive(\delta \te \circ \tau^{-1})$
	\item $A'(\delta \te \circ \tau^{-1})$
\end{itemize} 

These maps are:

\begin{itemize}
	\item $\delta \te_k \mapsto |\det(DT_k)|$
	\item $\delta \te_k \mapsto A_{T_k}$
\end{itemize}

Therefore:

\begin{itemize}
	\item $|\det(DT_k)|-1 = |\det(DT_k)|-1 - t_k\dive(\delta \te \circ \tau^{-1})+t_k\dive(\delta \te \circ \tau^{-1}) = o^1_k + t_k\dive(\delta \te \circ \tau^{-1})$
	\item $A_{T_k}-I = A_{T_k}-I - t_k A'(\delta \te \circ \tau^{-1}) + t_k A'(\delta \te \circ \tau^{-1}) = o_k^2 + t_k A'(\delta \te \circ \tau^{-1}) $
\end{itemize}

where $o_1^k \in L^\infty(\mR^n;\mR)$ and $o^2_k \in L^\infty(\mR^n;\mR^{n\times n})$ being higher order terms, in $L^\infty$ and with respect to $t_k$.

We can then write $(G(k,u^0,\pi^k)-G(0,u^0,\pi^k))/t_k = a_k + o_k$.

Here:

\begin{align*}
a_k :=\\
\frac{1}{2}\int_I \int_{T(U_r)}|v^\tau-w^\tau|^2\dive(\delta \te \circ \tau^{-1})+\\
\int_I ( w_t^\tau \dive(\delta \te \circ \tau^{-1}), q^k)_{H_T}+ (A'(\delta \te \circ \tau^{-1}) \nabla w^\tau, \nabla q^k)_{H_T}+\\
\int_I (v_t^\tau \dive(\delta \te \circ \tau^{-1}),p^k )_{H_T} + (A'(\delta \te \circ \tau^{-1})  \nabla v^\tau, \nabla p^k)_{H_T} \\
\end{align*}

Thanks to the assumed weak convergence, $a_k\rightarrow J'(\tau)[\delta \te]$.

So, we still have to show that:

\begin{align*}
o_k:=\\
\frac{1}{2}\int_I \int_{T(U_r)}|v^\tau-w^\tau|^2 o^1_k t_k^{-1}+\\
\int_I ( w_t^\tau o^1_kt_k^{-1}, q^k)_{H_T}+ (t_k^{-1}o^2_k\nabla w^\tau, \nabla q^k)_{H_T}+\\
\int_I (v_t^\tau o^1_k t_k^{-1},p^k )_{H_T} + (t_k^{-1}o^2_k \nabla v^\tau, \nabla p^k)_{H_T} 
\end{align*}

goes to zero. This is true because again we can write:

\begin{align*}
|o_k| \leq\\ \norm{ o^1_kt_k^{-1}}_{L^\infty(T(U_r))}\left (\int_I( \norm{v_t^\tau}_{H^T}\norm{p^k}_{H_T}+ \norm{v_t^\tau}_{H_T}\norm{p^k}_{H_T})+\frac{1}{2}\norm{v^\tau-w^\tau}^2_{L^2(I,H_T)}\right )+\\
\left(\sum_{ij} \norm{((t_k^{-1}o^2_k)_{ij}}_{L^\infty(T(U_r))}\right )\left (
\int_I \norm{\nabla v^\tau}_{H_T} \norm{\nabla p^k}_{H_T}+ \int_I \norm{\nabla w^\tau}_{H_T} \norm{\nabla q^k}_{H_T}\right )
\end{align*}

which goes to $0$, thanks to the boundedness of the averaged adjoint states, which stems from their weak convergence.

We assumed $p^k \weakc p^0$ in $Q(I,W_T)$ and $q^k \weakc q^0$ in $Q(I,V_T)$. We now prove these claims.

\underline{Weak convergence of states}

We show at first  $v_0^k \weakc v^0_0$ in $Q(I,W_T)$ and $w^k \weakc w^0$ in $Q(I,V_T)$.

We do this by showing a bound, uniform in $k$.

To do this, recall that $V_0^k:=v_0^k\circ T_k^{-1}$ and $W^k:=w^k\circ T_k^{-1}$ satisfy, as seen in \ref{thm:eq_pde}:

\begin{align*}
W^k \in Q_0(I, V_{T_k\circ T}), V_0^k \in Q_0(I,W_{T_k\circ T}) \\
\int_I  (W^k_t , Q)_{H_{T_k\circ T}}+ (\nabla W^k, \nabla Q)_{H_{T_k\circ T}} = \int_I(g,\tr_{{T_k\circ T}(U)} Q)_{L^2(\Gamma_f)}, \quad \forall Q \in Q^0(I, V_{T_k\circ T}) \\
\int_I (V^k_{0t},P)_{H_{T_k\circ T}} + (\nabla V_0^k, \nabla P)_{H_{T_k\circ T}}= -\int_I(U_k',P)_{H_{T_k\circ T}}+(U^k, \nabla P)_{H_{T_k\circ T}}, \quad \forall P \in Q^0(I, W_{T_k\circ T})
\end{align*}

where $U^k:=\bar{u}\circ T_k^{-1}$, where we used that pullbacks and time derivatives commute, see \cref{prop:change_boch}.

Thanks to \cref{prop:diri_wp} and \cref{prop:wp_neu} we obtain the stability estimates:

\begin{align*}
\norm{V^k}^2_{C([0;T],H_{T_k\circ T})}+\norm{V^k}_{L^2(I,H_{T_k\circ T})}^2+ \\\norm{\nabla V^k}_{L^2(I,H_{T_k\circ T})}^2 + \norm{(V^k)_t}^2_{L^2(I,H_{T_k\circ T})}\leq\\ C\norm{U^k}_{H^1(I,V_{T_k\circ T})}^2\\
\norm{W^k}^2_{C([0;T],H_{T_k\circ T})}+\norm{W^k}_{L^2(I,H_{T_k\circ T})}^2+ \norm{\nabla W^k}_{L^2(I,H_{T_k\circ T})}^2 +\\ \norm{W^k_t}^2_{L^2(I,H_{T_k\circ T})}\leq \\C\norm{g}_{H^1(I,L^2(\Gamma_f))}^2
\end{align*}


where $C$ is independent of $k$.

%But $\norm{U^k}_{H^1(I,V_{T_k\circ T})}^2 = \norm{\bar{u}\circ T_k^{-1}}_{H^1(I,V_{T_k\circ T})}^2\leq N_k\norm{\bar{u}}_{H^1(I,V_{ T})}^2$ by \cref{lemma:bochner_Hk_map}, where $N_k:=\norm{\circ T_k}_{L(V_{T_k\circ T},V_{T}))}$.

Now, consider \ref{thm:change}. It says that for almost every time:

\begin{align*}
\norm{U^k}_{V_{T_k\circ T}}\leq\\ \left ( 1+\norm{\det DT_k}_{L^\infty(\mR^n)}\right)^{1/2} \norm{(DT_k)^{-1}}_{L^\infty(\mR^n;\mR^{n\times n})}\norm{\bar{u}}_{H^1(T(U_r));\mR^n)}
\end{align*}

and the same goes for the first derivative. 


%The term $( \norm{\det DT_k}_{L^\infty(\mR^n)})^{1/2}$ is bounded with respect to $k$ as we already saw, and with analogous reasonings, so is $ \norm{(DT_k)^{-1}}$, thanks to the Frechet differentiability of $\delta \te_k \rightarrow  (DT_k)^{-1}$ established in lemme 4.3, page IV.8 \cite{murat}.

This bound is uniform on $k$ because of the continuity of the bound, with respect to $k$, as seen in 4.12, page IV.6, \cite{murat}.

We conclude that $\norm{U^k}_{H^1(I,V_{T_k\circ T})}^2$ is bounded and we thus have that $W^k \in Q_0(I, V_{T_k\circ T}), V_0^k \in Q_0(I,W_{T_k\circ T})$ are bounded.

Now, for almost all times, using 4.11, page IV.6 of \cite{murat}, we obtain that, for instance:

\begin{align*}
\norm{v_0^k}_{V_{ T}}\leq \\\left ( 1+\norm{\det (DT_k)^{-1}}_{L^\infty(\mR^n)}\right)^{1/2} \norm{DT_k}_{L^\infty(\mR^n;\mR^{n\times n})}\norm{V_0^k}_{H^1(T_K(T(U_r)))}
\end{align*}

where we remember that $H^1_0$ was chosen to be normed with the full $H^1$ norm.

The same goes for $w^k$ and the first derivatives in time, yielding that $w^k \in Q_0(I, V_{ T}), v_0^k \in Q_0(I,W_{T})$ are bounded.

We thus have $w^k\weakc w^? \in Q_0(I, V_{ T})$, $v_0^k \weakc v_0^? \in Q_0(I,W_{T})$, in the weak topologies of, respectively, $Q(I, V_{ T}), Q(I,W_{T})$, and modulo subsequences. The initial values are preserved becase $Q_0$ is closed and convex in the Hilbert space $Q$ (see \ref{prop:Q}). The closedness follows from the fact that the embedding into continuous function is linear bounded, and evaluation at $0$ is linear bounded from continuous functions.

We now prove that $w^?=w^0$, $v^?=v^0$, and this will yield the weak convergence of the whole sequence.

To prove e.g. that $v^?=v^0$, let us look at the weak formulations of $v_0^k$:

\begin{align*}
\int_I (v_{0t}^k,p |\det(DT_k)|)_{H_T} + (A_{T_k} \nabla v_0^k, \nabla p)_{H_T}+(\bar{u}',p|\det(DT_k)|)_{H_T}+(A_{T_k} \nabla \bar{u} , \nabla p)_{H_T} = 0
\end{align*}

for all $p \in Q_0(I,W_{T})$.

Let's analyize the first term, which is $\int_I (v_{0t}^k,p |\det(DT_k)|)_{H_T} ) =(v_{0t}^k,p |\det(DT_k)|)_{L^2(I,H_T)}$.

We can write:

\begin{align*}
(v_{0t}^k,p |\det(DT_k)|)_{L^2(I,H_T)} =\\ (v_{0t}^k,p )_{L^2(I,H_T)}- (v_{0t}^k,p )_{L^2(I,H_T)} + (v_{0t}^k,p |\det(DT_k)|)_{L^2(I,H_T)} =\\ (v_{0t}^k,p )_{L^2(I,H_T)} + (v_{0t}^k,p (|\det(DT_k)|-1))_{L^2(I,H_T)}
\end{align*}

Because $p \in  Q(I,W_{T})$, the first term converges to $(v_{0t}^?,p )_{L^2(I,H_T)}$, see \cref{prop:Q} for details on why we can write the time derivative of the limit. The other term can be estimated as follows:

\begin{align*}
|(v_{0t}^k,p (|\det(DT_k)|-1))_{L^2(I,H_T)}|\leq\\\norm{v_{0t}^k}_{L^2(I,H_T)}\norm{p (|\det(DT_k)|-1))}_{L^2(I,H_T)}\leq\norm{v_{0t}^k}_{L^2(I,H_T)}\norm{p }_{L^2(I,H_T)} \norm{|\det(DT_k)|-1}_{L^\infty}
\end{align*}

Where the first term in the product is bounded by the weak convergence property, and the last one goes to $0$ by continuity, see again 4.12, page IV.6 of \cite{murat}.

In a similar fashion for the other pieces, and by passing to the limit:
\begin{align*}
\int_I (v_{0t}^?,p)_{H_T} + (\nabla v_0^?, \nabla p)_{H_T}+(\bar{u}',p)_{H_T}+(\nabla \bar{u} , \nabla p)_{H_T} = 0
\end{align*}

By uniqueness, $v^?=v^0$.

\underline{Weak convergence of averages adjoint states}

So, $v_0^k \weakc v_0^0, w^k\weakc w^0$ in the sense of the $Q(I,W_T)$ and $Q(I,V_T)$ weak convergence.

We now claim that $p^k \weakc p^0, q^k\weakc q^0$, in a similar style as before. To do so, remember that $P^k:=p^k\circ T_k^{-1}$ and $Q^k:=q^k\circ T_k^{-1}$ solve:

\begin{align*}
-Q^k_t-\Delta Q^k =\frac{v_0^k-w^k+v_0^0-w^0}{2}\circ T_k^{-1}+\bar{u}_T\circ T_k^{-1} \\
Q^k(T)=0\\
\partial_\nu Q^k = 0 \text{ on } \Sigma_f\\
Q^k = 0 \text{ on } \Sigma_m
\end{align*}

and

\begin{align*}
-P^k_t-\Delta P^k =-\frac{v_0^k-w^k+v_0^0-w^0}{2}\circ T_k^{-1}-\bar{u}_T\circ T_k^{-1} \\
P^k(T)=0\\
P^k = 0 \text{ on } \Sigma_m \sqcup \Sigma_f
\end{align*}

By \cref{thm:const_track}, we will obtain a bound in $Q$ of the tranported averaged adjoints as soon as we have a bound on $\ds \frac{v_0^k-w^k+v_0^0-w^0}{2}\circ T_k^{-1}$ in the $L^2(I,H)$ norm, and of $U^k:=\bar{u}_T\circ T_k^{-1}$. The latter was proven above.

So, by \cref{thm:change} and 4.12 of \cite{murat} at page IV.6, it suffices to have an $L^2(I,H)$ bound on $\ds \frac{v_0^k+w^k+v_0^0+w^0}{2}\circ T_k^{-1}\circ T_k = \frac{v_0^k+w^k+v_0^0+w^0}{2}$ which we have, since we just proved that $v_0^k, w^k$ are weakly convergent in e.g. $L^2(I,H)$.

We conclude that $Q^k,P^k$ are bounded in the $Q(I,V_{T_k\circ T})$ and $Q(I,W_{T_k\circ T})$ sense.

But what we want is a bound on $q^k, p^k$ in the $Q(I,V_{ T})$ and $Q(I,W_{ T})$ sense. This can be accomplished in exactly the same way as before.

We conclude that there exist $q^?, p^?$ in $Q^0(I,V_{ T})$, $Q^0(I,W_{ T})$, that are, modulo subsequences, the weak limits of $q^k, p^k$.

To show e.g. $q^?=q^0$ and conclude the convergence of the full sequence, we analyze the weak formulation of $q^k$, which reads, after going to the moving domain and applying integration by parts in time (see \ref{prop:Q}):

\begin{align*}
-\int_I \left (\frac{((v_0^k+\bar{u}_T - w^k)+(v_0^0+\bar{u}_T - w^0)}{2}\right )|\det(DT_k)|,\delta w)_{H_T}+\\
-\int_I (  q^k_t ,   \delta w |\det(DT_k)|)_{H_T}+ (A_{T_k}\nabla \delta w, \nabla q^k)_{H_T}
\end{align*}

for all $\delta w \in Q_0(I,V_{ T})$.

We show the convergence of e.g. the member: $\int_I( v_0^k|\det(DT_k)|,\delta w)_{H_T}$. By splitting the scalar product as we saw above, we are left with checking that $\int_I (v_0^k,\delta w)_{H_T}\rightarrow  \int_I (v_0^0,\delta w)_{H_T}$, which is true, since we proved that $v_0^k \weakc v_0^0$ in $Q(I,W_T)$.

We conclude upon convergence that:

\begin{align*}
-\int_I \left (\frac{((v_0^0+\bar{u}_T - w^0)+(v_0^0+\bar{u}_T - w^0)}{2}\right ),\delta w)_{H_T}+\\
-\int_I (  q^?_t ,   \delta w )_{H_T}+ (\nabla \delta w, \nabla q^?)_{H_T}
\end{align*}

which is satisfied also by $q^0$, therefore $q^? = q^0$ and we have weak convergence of the entire sequence.

\end{mproof}

\appendix

\chapter{Functional spaces}
\section{Sobolev spaces}

\begin{thm}[Integration by parts]
\label{thm:ibp}
Let $\Omega$ be a bounded Lipschitz domain. Let $1<p<\infty$ and $f,g \in W^{1,p}(\Omega), W^{1,q}(\Omega)$, $q=p'$. Then:

$$\int_\Omega f \partial_i g = -\int_\Omega g \partial_i f+\int_{\partial \Omega} \tr u \nu_i d\pazocal{H}^{n-1}$$
\end{thm}
\begin{mproof}

This follows from \cite{leoni}, theorem 18.1 at page 592, where $g$ needs to be $C^1_c(\mR^n)$. But \cite{adams}, theorem 3.18 at page 54, says that (thanks to the smoothness of the boundary) the set of the restriction of such functions is dense in $W^{1,q}(\Omega)$, so that we can conclude by a density argument \textcolor{red}{\href{https://www.math.kit.edu/iana2/lehre/sobolevspaces2021s/media/lec15.pdf}{developed here}}.

\end{mproof}

\begin{lemma}
$f \in L^\infty(\Omega; \mR^N) \iff f_i \in L^\infty$, and two equivalent norms are $\norm{f}_a:=\norm{|f|}_\infty$, $\norm{f}_b:=\max_i\norm{f_i}_\infty$, for $|\cdot |$ any finite dimensional norm.
\end{lemma}
\begin{mproof}

We choose $|\cdot |=|\cdot|_1$.

Consider $f_n \in X_a = \{[f], f: \Omega \rightarrow \mR^n \text{ measurable }, \norm{f}_a\}$, Cauchy. Then every component is Cauchy in the scalar $L^\infty$, so that $f_n^i \rightarrow f^i $ in $L^\infty$. The limit $f$ is in $X_a$ because the functions $|f_i|$ are essentially bounded, and so is $|f|$. 

Then $\norm{f_n-f}_a\leq \norm{f_n-f_m}_a+\sum_i\norm{f_m^i-f^i}_\infty$ for all $n,m$. Choose $m\geq n$ with $\norm{f_m^i-f^i}\leq 1/(Nn)$ and conclude $X_a$ is Banach.

We know from \cite{leoni}, theorem B.88 at page 671, and page 669, we know that $X_b = \{[f], f: \Omega \rightarrow \mR^n \text{ measurable }, \norm{f}_b\}$ is Banach.

Moreover $X_a=X_b$ as sets, so that the thesis follows.


%We have that $\norm{f}_\infty = \sup_{\Omega\setminus X}f$ for all $X$ null sets on which $f \leq \norm{f}_\infty$ on $\Omega \setminus X$.
%
%In fact, that  $\norm{f}_\infty \geq \sup_{\Omega\setminus X}f$ is clear, whereas suppose there is $\epsilon >0$ with $\norm{f}_\infty -\epsilon \geq f$ on $\Omega \setminus X$. Then we contradict the definition of essential supremum and therefore $\norm{f}_\infty = \sup_{\Omega\setminus X}f$.
%
%Next up: for some $X,Y \subseteq \Omega$ of measure $0$, $f,g\leq \norm{f}_\infty, \norm{g}_\infty$ in $\Omega\setminus X, \Omega \setminus Y$, therefore, for some $N$ of null measure, $f,g\leq \norm{f}_\infty, \norm{g}_\infty$ in $\Omega \setminus N$. Therefore, $\norm{f}_\infty = \sup_{\Omega \setminus N}f, \norm{g}_\infty = \sup_{\Omega \setminus N}g$. 
%
%Therefore $ \max(\norm{f}_\infty, \norm{g}_\infty) = \max(\sup_{\Omega \setminus N}f,  \sup_{\Omega \setminus N}g) = \sup_{\Omega\setminus N}\max(f,g)$.
%
%So, $N$ is a null set and $\max(f,g)\leq \max(\norm{f}_\infty, \norm{g}_\infty)$, which by the reasoning of before.
%
%

\end{mproof}

\begin{prop}[Characterization of $W^{1,\infty}$]
\label{prop:lip}
Let $\Omega$ be a bounded Lipschitz domain, or $\mathbb{R}^n$. Then $W^{1,\infty}(\Omega) = C^{0,1}\cap L^\infty(\Omega)$.

This means that $u\in W^{1,\infty}(\Omega)$ if and only if $u$ has a (unique) representative that is bounded, Lipschitz continuous. Weak and classical derivatives coincide a.e.
\end{prop}
\begin{mproof}

%\underline{ACL characterization}
%
%Let $\Omega$ be just a domain.
%
%Consider a line $\gamma$ intersecting $\Omega$. Then, the intersection is a disjoint union of open segments $\gamma_i$. We say $u\in AC_\gamma(\Omega)$ if and only if $u$ is $AC[a,b]$ for every $[a,b]\cc\eta_i$ for all $i$, for almost all lines $\eta$ parallel to $\gamma$ and intersecting $\Omega$.
%
%We say $u \in ACL(\Omega)$ if $u$ is $AC_{e_i}(\Omega)$ for all coordinate axis $e_i$. We also call $BL(\Omega)$ the set of $u$ with a representative $\tilde{u}$, with $\tilde{u}\in ACL(\Omega), \tilde{u}, \nabla \tilde{u} \in L^\infty$, $\nabla u$ being the classical (a.e.) gradient.
%
%We have the following result (see \cite{kufner}, page 276, theorem 5.6.5.): $u \in W^{1,\infty}(\Omega)$ if and only if $u$ has a representative $\tilde{u} \in BL(\Omega)$. Classical and weak derivatives then coincide a.e..

\underline{Extension}

In the case $\Omega$ is bounded Lipschitz, then $\Omega$ is an extension domain for $W^{1,\infty}(\Omega)$, meaning that there is $E: W^{1,\infty}(\Omega)\rightarrow W^{1,\infty}(\mathbb{R}^n)$ linear bounded with $Eu=u$ a.e. on $\Omega$ (see \cite{leoni}, theorem 13.17 at page 425, 13.13 at page 424, and definition 9.57 at page 273).

\underline{The proof}

Let $u \in  W^{1,\infty}(\Omega)$. By \cite{leoni}, 11.50 at page 339, because $\Omega$ is an extension domain, we obtain that $u$ has a representative $\bar{u}$ that is bounded Lipschitz. Let $\phi \in C_c^\infty(\Omega)$. By The Kirszbraun theorem (see e.g. \cite{kirszbraun}), we can extend $\bar{u}$ to a Lipschitz function $e$ on $\mathbb{R}^n$. Then, for a large enough cube $Q$ containing $\Omega$, $\ds \int_\Omega\bar{u}\partial_i\phi = \int_Q e \partial_i \phi  = -\int_Q \partial_i e \phi $, by Fubini's theorem and integration by parts for AC functions.

Because $e=\bar{u}$	on $\Omega$, we conclude $\ds \int_\Omega\bar{u}\partial_i\phi =-\int_{\Omega} \partial_i \bar{u} \phi $, so that $\nabla \bar{u} = \nabla u$ almost everywhere.

Conversely, let $u$ be bounded Lipschitz. The above reasoning shows that $u$ has essentially bounded weak derivatives equal to the a.e. classical derivatives.

\end{mproof}

\begin{cor}[$W^{k,\infty}=C^{k,1}_B$]
\label{prop:lipk}
For a bounded Lipschitz domain $\Omega$, or for $\Omega = \mR^n$, then $W^{k,\infty}=C^{k,1}_B$ ($C^{k,1}$ bounded functions with bounded classical derivatives up to order $k+1$).
\end{cor}
\begin{mproof}

We have already proved the case $k=1$. We prove, for instance, the case $k=2$. Then, $u \in W^{k,2} \implies u, \partial_i u \in W^{k,1}$ (\cite{leoni}, 11.7 at page 321), so that by \cref{prop:lip}, we find bounded Lipschitz $h, g_i$ with $u=h$ a.e., $\partial_i u = \partial_i h$ a.e., $g_i = \partial_i u$ a.e..

Therefore $h$ is continuous, with continuous weak derivatives $g_i$, which implies that $h \in C^1(\Omega)$ (see \textcolor{red}{\href{https://math.stackexchange.com/questions/497708/is-a-continuous-function-with-continuous-weak-derivatives-of-class-c1}{here} and \href{https://math.stackexchange.com/questions/1787716/u-continuous-and-the-weak-derivative-du-continuous-rightarrow-u-in-c1}{here}}).

Now, $\partial_i h =g_i$ a.e., so everywhere, so that:

\begin{itemize}
\item $h$ is bounded Lipschitz and $C^1$
\item $\partial_i h$ are bounded Lipschitz
\end{itemize}

%Therefore, for $\phi \in C_c^\infty(\Omega)$ we get $\ds -\int_\Omega u \phi_{,i}\ds = \int_\Omega \partial_i u\phi =  \int_\Omega g_i \phi$.
%
%Pick a cube $Q = [-M,M]^n$, $\Omega \cc Q$. So, upon extending in a bounded Lipschitz manner $g_i$ to $G_i$ with $G_i(\partial Q)=0$ (see \cite{kirszbraun} and multiply by a suitable cut-off function), and $\phi$ to $0$, we get $\ds ...=\int_\Omega G_i \phi$. Call $h_i^x(t)=G_i(...,x_{i-1},t,x_{i+1},...)$ which is $AC[-M,M]$. Because $\eta_i^x(s):= \int_{-M}^s h_i^x(z)dz$ is $AC[-M,M]$ too because $h_i^x$ is bounded, and $h_i^x(t) =  \eta_i^{x'}(t)$. Employing then Fubini:
%
%\begin{align*}
%...= \int_{-M}^M... \int_{-M}^M \int_{-M}^M G_i(...,x_{i-1},x_i,x_{i+1},...)\phi dx_i dx_1...dx_{i-1}dx_{i+1}...dx_n  \\
%= \int_{-M}^M... \int_{-M}^M \int_{-M}^M h_i^x(x_i) \phi dx_i dx_1...dx_{i-1}dx_{i+1}...dx_n \\
%= \int_{-M}^M... \int_{-M}^M \int_{-M}^M \eta_i^{x'}(x_i)  \phi dx_i dx_1...dx_{i-1}dx_{i+1}...dx_n \\
%= \int_{-M}^M... \int_{-M}^M \int_{-M}^M \eta_i^{x'}(x_i)  \phi dx_i dx_1...dx_{i-1}dx_{i+1}...dx_n \\
%= -\int_{-M}^M... \int_{-M}^M \int_{-M}^M \eta_i^x(x_i)  \phi_{,i} dx_i dx_1...dx_{i-1}dx_{i+1}...dx_n \\
%=- \int_{-M}^M... \int_{-M}^M \int_{-M}^M \int_{-M}^{x_i} G_i(...,x_{i-1},z,x_{i+1},...) dz\phi_{,i} dx_i dx_1...dx_{i-1}dx_{i+1}...dx_n \\
%= -\int_Q k_i\phi_{,i}
%\end{align*}
%
%where in the last passage we have used again Fubini and the fact that 
%
%$$k_i(x):=\int_{-M}^{x_i} G_i(...,x_{i-1},z,x_{i+1},...) dz$$
%
% is measurable and bounded on $Q$.

\end{mproof}

\section{Bochner spaces}

Here are some useful results about Bochner spaces.

\begin{prop}[Bochner integral and bounded operators]
\label{prop:bochner_bound}
Let $X,Y$ be separable Banach, let $T \in L(X,T)$ be a linear bounded operator. For $f \in L^1(I,X)$ define $Tf (t):= T(f(t))$. Then $Tf \in L^1(I,Y)$ with $T\int_I f = \int_I Tf$.
\end{prop}
\begin{mproof}

First of all, a clarification on the definition. What is really happening is that from the time equivalence class $f$, we select a $g$, and then $Tf(t):=T(g(t))$. $Tf$ is then the equivalence class of $t\mapsto Tf(t)$. The definition is well posed, because $g_1(t)=g_2(t)\implies T(g_1(t))=T(g_2(t))$.

Let $f_n$ be simple, $f_n\rightarrow f $ a.e., with $\lim_n \int_I f_n = \int_I f$ and $\norm{f_n}_X \leq C \norm{f}_X$ (see page 6, and corollary 2.7 at page 8 of \cite{kreuter}).

\underline{Measurability}

For almost all $t$, $T(f_n(t)) \rightarrow T(f(t))=Tf(t)$ in $Y$, so that $Tf$ is measurable (strongly).

\underline{Integrability}

By the assumptions, $\norm{Tf_n}\leq \norm{T}\norm{f_n}\leq C\norm{f}\in L^1(I)$, so that by dominated convergence (corollary 2.6 of \cite{kreuter}) $Tf$ is integrable too. Thus $\int_T Tf = \lim_n \int_I  Tf_n = \lim_n T\int_I  f_n$, because $f_n$ is simple. And now, by the choice of $f_n$, $\int_T Tf = \lim_n T\int_I  f_n = T \lim_n \int_I  f_n = T \int_I f$.

\end{mproof}

\begin{prop}	[Derivations and bounded operators]
\label{lemma:bochner_Hk_map}
As before, let $X,Y$ be separable Banach, let $T \in L(X,T)$ be a linear bounded operator.

For $k\geq 0$, $f \in H^k(I,X)\implies Tf \in H^k(I,Y)$, with weak derivatives $\partial_{t^i}Tf = T\partial_{t^i}f$, $0\leq i \leq k$.

The map $f \mapsto Tf$, $H^k(I,X)\rightarrow H^k(I,Y)$ is linear, and bounded by $\norm{T}$.
\end{prop}
\begin{mproof}
The case $k=0$ is proved above.

We prove now that $\partial_{t^i}Tf = T\partial_{t^i}f$ for $i=1$. Note that $T\partial_t f \in L^2(I,Y)$, which qualifies as weak derivative.

In fact, for $\phi \in C_c^\infty(I)$, we have $\int_I \phi T\partial_i f = \int_I T(\phi\partial_t f) = T \int_I\phi\partial_t f = -T\int_I\phi'f=-\int_I\phi'Tf$.

Higher weak derivatives are treated analogously and the rest of the claims follow from the time stationarity of $T$ and by $\norm{\partial_{t^i}Tf}=\norm{T\partial_{t^i}f}\leq \norm{T}\norm{\partial_{t^i}f}$.

\end{mproof}

\begin{prop}[Continuous representatives]
\label{prop:cts_repr}
Let $X$ be separable Banach. $f \in L^1(I,X)$ has at most a continuous representative on $[0,T]$.
\end{prop}
\begin{mproof}
Assume there exists two such continuous representatives, so that we get a function $\delta: [0,T] \rightarrow X$ that is zero almost everywhere and continuous. Hence, $[0,T] \ni t \mapsto \norm{\delta(t)}$ is continuous in $\mR$ and zero a.e., so that it must be zero everywhere.
\end{mproof}

%\begin{prop}[Weak derivatives and Gelfand triples]
%\label{prop:sob_implies_W}
%Let $H\subseteq V$ be separable Hilbert spaces, $H$ densely embedded in $V$.
%
%Then $y \in L^2(I,V)\cap H^1(I,H) \subseteq W(I,V)$.
%\end{prop}
%\begin{mproof}
%We note that $V$ is Hilbert separable, so, reflexive Banach and separable, so that $V^*$ is separable too. Then, call $i$ the embedding $H \emb V^*$.
%
%For $\phi \in C^\infty_c(I)$ we obtain $\int_I i(y_t)\phi=\int_I i(y_t\phi)= i \left ( \int_I y_t\phi \right )$ by \cref{prop:bochner_bound}. Therefore $\int_I i(y_t)\phi=-i \left ( \int_I y\phi' \right ) =\int_I i(y)\phi'  $ , proving that $y \in W(I,V)$.
%
%\end{mproof}

We now check that a vector valued test function has weak derivatives of all orders.

\begin{prop}[Weak derivatives of test functions]
\label{prop:weak_class}
Let $\phi \in C^1([0,T],X)$, for $X$ separable Banach. It means that the limit of the difference quotients exists for all points of $I$, that $t\mapsto \phi(t), \phi'(t)$ are continuous, and that they can be continuously extended to $[0,T]$.

Then these classical derivatives coincide a.e. with the weak derivatives of $u$.

\end{prop}
\begin{mproof}


We rely on proposition 3.8 of \cite{kreuter} at page 26.

\underline{Absolute continuity}

Consider $\epsilon >0$. Divide $[a,b] \cc (0,T)$ into a uniform partition $t_i$. By theorem 6 at page 146 of \cite{mvt}, we get that $\ds \norm{\phi(t_i)-\phi(t_{i-1})}_X \leq (t_i-t_{i-1})\norm{\phi(\xi_i)}_X\leq (b-a) \norm{\phi'}_{\infty}/n$, and by choosing $n$ small enough, we conclude that $\phi$ is (locally) absolutely continuous.

\underline{Weak derivative}

Therefore, $\phi$ is locally $AC$, differentialble everywhere and $\phi'$ is bounded, so that $\phi \in H^1(I,X)$ and weak and classical derivatives coincide. 

\end{mproof}

And now, introduce a time dependent version of the trace operator which is useful for our computations.

\begin{defn}[Time dependent trace]
Let $\Omega$ be a bounded Lipschitz domain. For $k\geq 0$ we define $\tr: H^k(I,H^1(\Omega))\rightarrow H^k(I, H^{1/2} (\partial \Omega))$ by $\tr(u)(t):=\tr(u(t))$
\end{defn}

Below are some properties of this operator.

\begin{prop}[Properties of trace operator]
\label{prop:trace}
The trace operator just defined:
\begin{enumerate}
\item is well posed
\item is linear bounded
\item admits a linear bounded right inverse, for instance, $E(g)(t):=E(g(t))$ (for $E$ a right inverse of the static trace)
\item $\tr$ and $E$, in the case of $k \in \mathbb{N}_0$, coincide (in the time a.e. sense) for the case $l\geq k$
\item for $k\geq 1$, $\tr u(0)=0 \iff u(0)=0$ (in the sense of continuous representatives)
\item it coincides with the trace treated for instance in \cite{lions}
\end{enumerate}
\end{prop}
\begin{mproof}

\underline{Proof of the proposition}

We recall that the trace operator is bounded surjective onto $H^{1/2}(\partial \Omega)$, with a right inverse $E$ (see theorem 3.37 at page 102 of \cite{mclean}).

The first three points are consequences of this fact and of \cref{prop:bochner_bound}.

The fourth property follows by the definition of $\tr, E$ and the fact that $H^l\subseteq H^k$, for $k\leq l$.

Let now $k\geq 1$. We know that $H^1, H^{1/2}$ are separable and Banach (the latter is separable because the continuos image of $H^1$ separable, and Banach (see \cite{grisvard}, page 20). Therefore, by \cite{evans}, theorem 2 of page 286, we obtain the embeddings $H^k(I,H^1)\emb C([0,T],H^1)$ and the same goes for $H^k(I,H^{1/2})$. The embedding is $U$, the unique continuous representative of a certain time equivalence class (\cref{prop:cts_repr}). We also introduce brackets to indicate equivalence classes in time, so, $u = [Uu]$.

We want to prove $(Uu)(0)=0 \iff U(\tr u)(0)=0$. But we have $[t \mapsto U(\tr u )(t)]=tru:=[t \mapsto \tr((Uu)(t))]$. So, $U(\tr u )(t)=\tr((Uu)(t))$ for all $t\in[0,T]$ by continuity. 

For the last point, let $k=0$. We have:

\begin{enumerate}
\item $H^1(\Omega)\cap C^1(\overline{\Omega})$ is dense in $H^1(\Omega)$ (see \cite{adams}, theorem 3.18 at page 54, where being $\Omega$ bounded Lipschitz is important)
\item functions $\sum_{i\leq m} \phi_i(t)f_i$ for $\phi_i \in C_c^\infty(I), f_i \in H^1(\Omega)\cap C^1(\overline{\Omega})$ are dense in $L^2(I,H^1)$ (see \cite{hinze}, page 39, lemma 1.9)
\end{enumerate}

It follows by the third point that $C^1(\overline{\Omega\times I})$ is dense in $L^2(I,H^1)$, so that $u\mapsto u|_{I\times \partial \Omega}$ admits a unique extension by continuity to $L^2(I,H^1)$, so that this definition of trace coincides with the one from the literature in the case of the space $H^{1,0}:=L^2(I,H^1)$ (see \cite{lions}, theorem 4.1), we expand this argument below.

\underline{Proof of leftover facts}

We call $C^k(\overline{\Omega}):=\{u \in C^k(\Omega) \text{ with }\partial_\alpha f \text{ extendable by continuity to } \overline{\Omega} \}$. 

Consider $u(x,t):=\phi(t)v(x)$, for $\phi \in C^1([0,T]), v \in C^1(\overline{\Omega})$. Then, it has partial derivatives $u_t = \phi_t v, u_i = \phi u_i$. $u$ and all its partial derivatives are continuous on $I\times \Omega$, meaning that $u \in C^1(\Omega \times I)$.

Moreover, $u, u_i, u_t \in C([0,T], C(\overline{\Omega}))$. We claim $ C([0,T], C(\overline{\Omega})) = C(\overline{\Omega\times I})$. In fact, one direction is trivial, and so, let $f \in C([0,T], C(\overline{\Omega})) = C(\overline{\Omega})$. Fix $(t,x) \in \overline{\Omega\times I}$. Then, $|f(s,y)-f(t,x)|\leq |f(t,y)-f(t,x)|+|f(t,y)-f(s,y)|\leq  |f(t,y)-f(t,x)|+\norm{f(t, \cdot)-f(s,\cdot)}_{\infty}$. If now $s$ is close to $t$, and $y$ is close to $x$, then $|f(s,y)-f(t,x)|$ is small.

This shows $u, u_i, u_t \in C([0,T], C(\overline{\Omega})) \in C(\overline{\Omega\times I}) $, i.e. $u \in C^1(\overline{Q\times I})$.

To conclude, let $u \in L^2(I,H^1)$. Approximate $u$ by $u_k:=\sum_{i\leq m_k} \phi_i^k(t)f_i^k$ as in point 2, and approximate $f_i^k$ by suitable $g_i^k \in H^1(\Omega)\cap C^1(\overline{\Omega})$, to obtain $u_k:=\sum_{i\leq m_k} \phi_i^k(t)g_i^k$

Then $\norm{u-w_k}_{L^2(I,H^1)}\leq \norm{u_k-w_k}_{L^2(I,H^1)}+\norm{u_k-u}_{L^2(I,H^1)}$. We only need to estimate $ \norm{u_k-w_k}_{L^2(I,H^1)}\leq\ds  T \sum_{i\leq m_k} \norm{\phi_i^k}_\infty\norm{f_i^k-g_i^k}_{H^1}$. By the first point, $\norm{f_i^k-g_i^k}_{H^1}$ can be made as small as it is necessary to conclude.

\underline{Last remarks}

Again with reference to \cite{lions}, consider the anisotropic spaces $H^{r,s}:=L^2(I,H^r)\cap H^s(I,L^2)$. We restrict to the case $r = 1$, $s\geq 0$. Denote the traces $\tr_s$ defined in theorem 4.1, mapping $H^{1,s}(\Omega \times I)\rightarrow H^{1/2, s/2}(\partial \Omega \times I)$. For $\partial \Omega$ Lipschitz this theorem is still valid, as $1/2\leq 1$, see the discussion above lemma 2.4 in \cite{costabel}. As stated in \cite{lions}, $\tr_s$ is an extension of  $u\mapsto u|_{I\times \partial \Omega}$, defined on the dense suspace $C^\infty(\overline{Q\times I})$ of $H^{1,s}$ (that this space is dense can be proved as in lemma 2.22 of \cite{costabel}). So, let $C^\infty(\overline{Q\times I})\ni u_n \rightarrow_{H^{r,s}} u \in H^{1,s}$.
 
We have $\tr_s u_n = \tr_0 u_n$. Then, $u_n \rightarrow_{H^{1,s}} u$, $u_n \rightarrow_{H^{1,0}} u$, so that $\tr_s u_n \rightarrow_{H^{1/2,s/2}}\tr_s u$ (hence $\tr_0 u_n \rightarrow_{H^{1/2,0}}\tr_s u$) and $\tr_0 u_n  \rightarrow_{H^{1/2,0}} \tr_\sigma u$.

Thus $\tr_0 u = \tr_s u$.

Using what we derived before, we can conclude the characterization of the traces in the anisotropic settign define 

\end{mproof}

And now some sanity checks in the case of Gelfand triples. 

\begin{prop}[Sanity checks for Gelfand triples]
\label{prop:sanity}

Consider the following Gelfand triples (the diagram commutes):

\[\begin{tikzcd}
	{V} &&& {V^*} \\
	& H & {H^*} \\
	{W} &&& {W^*}
	\arrow["c", hook, from=3-1, to=1-1]
	\arrow["a", hook, from=1-1, to=2-2]
	\arrow["b"', hook, from=3-1, to=2-2]
	\arrow["{a^*}", hook, from=2-3, to=1-4]
	\arrow["r", from=2-2, to=2-3]
	\arrow["{b^*}"', hook, from=2-3, to=3-4]
	\arrow["{c^*}", hook, from=1-4, to=3-4]
\end{tikzcd}\]

Here $W\subseteq V \subseteq H$ are all separable Hilbert spaces, $a,b,c$ the trivial injections, $r$ the Riesz isomorphism of $H$. We denote by $i_V$ the Gelfand triple embedding $V\emb V^*$, so, $i_V=a^*ra$.

Then:

\begin{enumerate}
	\item $H^1(I,V)\subseteq W(I,V)$ with continuous embedding. The $W(I,V)$ derivative of $u \in H^1(I,V)$ is $i_V u_t$.
	\item for $u \in W(I,W)$ with $(i_W u)'\in L^2(I,H)$ (i.e. $(i_W u)_t = b^*r h$ for $h$ in $L^2(I,H)$) we obtain $u \in W(I,V)$ (i.e. $cu \in W(I,V)$), with derivative $(i_V cu)'=a^*r h$, so that also $(i_V cu)' \in L^2(I,H)$. It also holds $(i_V cu)'|_W = (i_W u)'$. $h$ is also the weak derivative $L^2(I,H)$ of $bu$.
	\item let $u, v \in W(I,V)$ with $u-v \in W$. Then $u-v \in W(I,W)$ with derivative $(i_W(u-v))'=(i_V u)'|_W-(i_V v)'|_W$.
\end{enumerate}

\end{prop}
\begin{mproof}

We use several times that time integrals and bounded linear static operators commute, see \cref{prop:bochner_bound}. $\phi$ denotes $\phi \in C^\infty_c(I)$.

\underline{First point}

We need to check that $a^* r a u \in H^1(I, V^*)$. This follows from \cref{lemma:bochner_Hk_map}, so that $(a^* r a u)_t = a^* r a u_t$.

\underline{Second point}

At first we claim that $h$ is a weak derivative of $bu \in L^2(I,H)$. In fact, $b^* r \int_I bu \phi' = \int_I (i_W u)\phi' = \ind{\mind{u \in W(I,W)}} =-\int_I(i_Wu)'\phi=-\int_I b^*r h\phi = b^*r (-\int_I h \phi)$. By density (definition of Gelfand triple), $b^*$ is injective, $r$ is too, and thus $\int_I bu \phi' = -\int_I h \phi$, which shows that $bu$ has weak derivative $h$, in the $H^1(I,H)$ sense.

And now $\int_I i_V c u \phi'= \int_I a^*racu \phi' = a^*r\int_I bu\phi' = \ind{by what we just proved}=-a^*r \int_I h \phi$, proving that $(i_V cu)'=a^*r h$.

Morevoer $(i_V cu)'|_W = c^*a^*r h=b^*rh=\ind{assumption} = (i_W u)'$.

\underline{Third point}

We check the derivative. We have $\int_i i_W(u-v)\phi' =\ind{\mind{u-v \in W\subseteq V}} = \int_I b^* r a (u-v) = c^*\int_I(i_Vu - i_Vv)\phi' = -\int c^*((i_V u)'- (i_V v)')\phi$.

\end{mproof}

\chapter{Parabolic equations}
\label{chap:parab_eq}

\section{Abstract theory}

\begin{ass}[Basic assumption for parabolic problems]
\label{ass:basic_par}

Let $V\subseteq H$ be real separable Hilbert spaces, $V$ dense in $H$. Then $H\hookrightarrow V^*$ is also dense, as stated in \cite{trol} at page 147. This embedding is $H \ni f \mapsto (f, \cdot )_H$. We thus obtain a Gelfand triple, and we have $W(I,V)\subseteq C(I,H)$.

Let $A:V\rightarrow V^* $ be linear bounded, $u \in W(I;V)$, $f \in L^2(I,V^*)$ and $u_0 \in H$.

We also assume that $\langle Av, v \rangle_{V^*,V}+ \lambda \HN{v}^2\geq \alpha \VN{v}^2$ for $\lambda \geq 0, \alpha >0$.
\end{ass}

We are interested in the following problem:

\begin{pb}[Abstract parabolic equation]
\label{eqn:general_parabolic}
\begin{align}
	u_t+Au=f \text{ in }V^* \text{ and for a.e. } t \in (0,T)\\
	u(0)=u_0
\end{align}
\end{pb}

\begin{thm}[Basic well posedness of \cref{eqn:general_parabolic}]
\label{thm:well_pos_parabolic}
Under \cref{ass:basic_par}, \cref{eqn:general_parabolic} has a unique solution $u$. Moreover $u$ satisfies the energy estimate:
\begin{equation}
	\label{eqn:en_est}
	\norm{u}_{W(I,V)} + \norm{u}_{C([0,T],H)}\leq c(\lambda, \alpha, \VSN{A}, T)(\HN{u_0}+\norm{f}_{L^2(I,V^*)})
\end{equation} 
\end{thm}
\begin{mproof}
See \cite{gilardi} at page 19, theorem 26.
\end{mproof}

We can also obtain additional regularity. Here are further assumptions to make this possible.

\begin{ass}[Assumptions for additional regularity]
\label{ass:reg_par}
We assume $u_0 \in V$, $f = f_1+f_2 \in L^2(I,H)+H^1(I,V^*)$. We also need $A$ to be symmetric (i.e. $\langle Au,v \rangle_{V^*,V} = \langle Av,u \rangle_{V^*,V}$).
\end{ass}

\begin{thm}[Regularity of time derivative]
\label{thm:reg_time}
Suppose \cref{ass:basic_par} and \cref{ass:reg_par}. Then $u_t \in L^2(I, H)$ with the estimate:
\begin{align}
	\norm{u}_{W(I,V)} + \norm{u}_{C(I,H)} + \norm{u_t}_{L^2(I,H)} \leq\\ c(\lambda, \alpha, \VSN{A}, T)(\VN{u_0}+\norm{f_1}_{L^2(I,H)} + \norm{f_2}_{H^1(I,V^*)})
\end{align}

That $u_t \in L^2(I, H)$ means precisely that there is $h \in L^2(I,H)$ with $a^* r h = (i_V u)'$, with the notation introduced in \cref{prop:sanity}.

\end{thm}
\begin{mproof}
We refer to page 26 of \cite{gilardi}, theorem 28, and only prove the necessary modifications.

\underline{Product rule for $A$}

We have 

\begin{align*}
\int_0^t \langle Au_n,u_n' \rangle_{V^*,V} = \sum_{k,l\leq n} \langle Aw^n_k,w^n_l \rangle_{V^*,V} \int_0^t g_k^n g_l^{n '} = \\
\sum_{k,l\leq n} \langle Aw^n_k,w^n_l \rangle_{V^*,V} \left (-\int_0^t g_k^{n'} g_l^n + g_k^n(t) g_l^n(t)-g_k^n(0) g_l^n(0) \right )
\end{align*}

By linearity at first and then symmetry we get:

\begin{align*}
	= \langle Au_n,u_n \rangle_{V^*,V}(t)-\langle Au_n,u_n' \rangle_{V^*,V}(0)-\int_0^t\langle Au_n',u_n \rangle_{V^*,V}=\\
	= \langle Au_n,u_n \rangle_{V^*,V}(t)-\langle Au_n,u_n' \rangle_{V^*,V}(0)-\int_0^t\langle Au_n,u_n' \rangle_{V^*,V}
\end{align*}

so that:

\begin{align*}
	\int_0^t \langle Au_n,u_n' \rangle_{V^*,V} = \frac{1}{2}\left( \langle Au_n,u_n \rangle_{V^*,V}(t)-\langle Au_n,u_n' \rangle_{V^*,V}(0) \right )
\end{align*}

\underline{Estimate for right hand side}

We have:

\begin{align*}
	\int_0^t \langle f_2,u_n' \rangle_{V^*,V} = \sum_{k\leq n}\int_0^tg_k^{n'}\langle f_2,w_k^n \rangle_{V^*,V}
\end{align*}

By the smoothness of $f_2$ we have that $t \mapsto \langle f_2(t),w_k^n \rangle_{V^*,V}$ is $H^1(0,T)$, in particular $AC[0,t]$, so that we can integrate by parts:

\begin{align*}
	 = - \sum_{k\leq n}\int_0^tg_k^{n}\langle f_2',w_k^n \rangle_{V^*,V} + \sum_{k\leq n} g_k^{n}(t)\langle f_2(t),w_k^n \rangle_{V^*,V} - \sum_{k\leq n} g_k^{n}(0)\langle f_2(0),w_k^n \rangle_{V^*,V} = \\
	 -\int_0^t \langle f_2',u_n \rangle_{V^*,V} + \langle f_2,u_n \rangle_{V^*,V}(t)-\langle f_2,u_n \rangle_{V^*,V}(0)
\end{align*}

Here we have used \cref{lemma:bochner_Hk_map} to take the derivative inside the bracket.

Note that by the smoothness of $f_2$, we can write, for instance, $\langle f_2,u_n \rangle_{V^*,V}(0) = \langle f_2(0),u_n \rangle_{V^*,V}$.

\tred{NB: here I need $f_2(0) \in V$ probably, but I'm not yet using the compatibility condition!}.

Going to the absolute values:

\begin{align*}
	\left | \int_0^t \langle f_2,u_n' \rangle_{V^*,V}\right | \leq 
	\int_0^T |\langle f_2',u_n \rangle_{V^*,V}|+ \VSN{f_2(t)}\VN{u_n(t)}+\VSN{f_2(0)}\VN{u_n(0)} \leq \\
	\frac{1}{2}\norm{f_2'}_{L^2(I,V^*)}^2 + \frac{1}{2}\norm{u_n}_{L^2(I,V)}^2 + \frac{\alpha}{4}\VN{u_n(t)}^2 +\\
	+ \frac{4}{\alpha}\norm{f_2}_{L^\infty(I,V^*)}^2+ \frac{1}{2}\norm{f_2}_{L^\infty(I,V^*)}^2+\frac{1}{2}\VN{u_{n0}}^2
\end{align*}

Now, $u_n$ converges weakly in $L^2(I,V)$ by estimate (59) of \cite{gilardi} and thus $\frac{1}{2}\norm{u_n}_{L^2(I,V)}$ is bounded. The term $\frac{\alpha}{4}\VN{u_n(t)}$ can be pulled to the left hand side, $u_{0n}$ is $V$ convergent hence bounded. Therefore as in \cite{gilardi} we are able to conclude that $u_n'$ is bounded in $L^2(I,H)$. We want to conclude $u_t \in L^2(I,H)$. We know for sure that $\langle u_m',w_j\rangle_{V^*,V}=\langle f-Au_m,w_j\rangle_{V^*,V}$, so that muliplication by a test function and integration yields $\int_I \langle u_m',w_j\phi\rangle_{V^*,V}=\int_I \langle f-Au_m,w_j\phi\rangle_{V^*,V}$. Because $u_m\rightharpoonup u$ in $L^2(I,V)$ we observe that, by \cref{prop:bochner_bound} applied on $A \in L(V,V^*)$, it holds $\int_I \langle u_m',w_j\phi\rangle_{V^*,V}\rightarrow \int_I \langle u',w_j\phi\rangle_{V^*,V}$.

What's more, is that $u_m' \rightharpoonup h$ in $L^2(I,H)$, so that $\int_I \langle h ,w_j\rangle_{V^*,V}\phi = \int_I \langle u',w_j\rangle_{V^*,V}\phi$. It means that for almost all $t$, $\langle h ,w_j\rangle_{V^*,V} = \langle u',w_j\rangle_{V^*,V}$. And now we can really say that $u' \in L^2(I,H)$, which even more precisely means $(i_Vu)' = a^* r h$ almost everywhere.

We also obtain that $u_t$ is bounded by $c(\alpha)(\norm{f_2}_{L^\infty(I,V^*)}+\norm{f_2}_{L^2(I,V^*)}+\VN{u_0}+\norm{u}_{L^2(I,V)})$.

Note that, by \cite{evans}, theorem 2 of page 286, we can estimate $\norm{f_2}_{L^\infty(I,V^*)}$ by $c(T)\norm{f_2}_{H^1(I,V^*)}$, so that the claim for the time derivative $u_t$ is proven.

\end{mproof}

For the case where $H=L^2$, $H^1\supseteq V\supseteq H^1_0$,  $f_2|_{H^1_0}=0$ we have even more regularity available.

\begin{thm}[Additional regularity]
\label{thm:par_reg}
Suppose \cref{ass:basic_par} and \cref{ass:reg_par}. 

Let additionally $H=L^2$, $H^1\supseteq V\supseteq H^1_0$,  $f_2|_{H^1_0}=0$. Then $Au|_{H^1_0}$ extends to $\overline{Au_{H^1_0}} \in L^2(I,H)$ with:
\begin{align}
	\norm{u}_{W(I,V)} + \norm{u}_{C([0,T],H)} + \norm{u_t}_{L^2(I,H)} +\norm{\overline{Au|_{H^1_0}}}_{L^2(I,H)}\leq\\ c(\lambda, \alpha, \VSN{A}, T)(\VN{u_0}+\norm{f_1}_{L^2(I,H)} + \norm{f_2}_{H^1(I,V^*)})
\end{align}

Moreover $u_t+\overline{Au_{H^1_0}}=f_1$ in $L^2(0,T,L^2)\cong L^2(Q)$ and $\overline{Au|_{H^1_0}}=Au$ on $H^1_0$.

\end{thm}
\begin{mproof}

For $v \in H^1_0$ we get $ \langle Au,v \rangle_{V^*,V} =  \langle f_1-u_t,v \rangle_{V^*,V} = ( f_1-u_t,v )_H$, for almost all $t \in (0,T)$. From here we conclude that $Au(t)$ extends for a.a. $t$ to an element of $H$ with $(\overline{Au}-f_1+u_t,v)_{L^2}=0$ for all $v \in H^1_0$, almost all $t$. By density, $\overline{Au}-f_1+u_t=0$ in $H$ for almost all $t$, so that $\overline{Au}=f_1-u_t$ in $L^2(0,T,L^2)\cong L^2(Q)$.

This isometric isomorphism is stated in \cite{trol}, page 144. 

\end{mproof}

For our applications we also need to track the constants more precisely, which is accomplished in the next proposition.

\begin{prop}[Tracking the costants]
\label{thm:const_track}
With \cref{ass:basic_par} there holds:
\begin{align}
\norm{u}^2_{C([0;T],H)}+\alpha\norm{u}_{L^2(I,V)}^2\leq \exp(2\lambda T)(\HN{u_0}^2+\alpha^{-1}\norm{f}^2_{L^2(I,V^*)})\\
\norm{u'}_{L^2(I,V^*)}\leq \norm{A}_{L(V,V^*)}\alpha^{-1/2}\sqrt{\exp(2\lambda T)}\HN{u_0} +\\\left (\norm{A}_{L(V,V^*)}\alpha^{-1}\sqrt{\exp(2\lambda T)}+1\right ) \norm{f}_{L^2(I,V^*)}
\end{align}

With additionally \cref{ass:reg_par} we obtain:

\begin{align}
C\norm{u'}^2_{L^2(I,H)}\leq 
(1+(1+C_0)\alpha^{-1})\norm{f_2}_{H^1(I,V^*)}^2+\\
(1+\norm{A}_{L(V,V^*)})\VN{u_{0}}^2+C_0\HN{u_0}^2+\\
\norm{f_1}_{L^2(I,H)}^2+C_0\alpha^{-1}\norm{f_1}^2_{L^2(I,V^*)}
\end{align}

with $C>0$ a number independent of the problem.

Here $C_0 = \ds 2^{-1}\max(1,\lambda)\max(1,\alpha^{-1})\exp(2\lambda T)$.

\end{prop}
\begin{mproof}

\underline{No regularity}

From page 21 of \cite{gilardi} we obtain that $\norm{u}^2_{C([0;T],H)}+\alpha\norm{u}_{L^2(I,V)}^2\leq \exp(2\lambda T)(\HN{u_0}^2+\alpha^{-1}\norm{f}^2_{L^2(I,V^*)})$.

In particular, $\sqrt{\alpha}\norm{u}_{L^2(I,V)}\leq 	\sqrt{\exp(2\lambda T)}(\HN{u_0}+\alpha^{-1/2}\norm{f}_{L^2(I,V^*)})$, or $\norm{u}_{L^2(I,V)}\leq 	\alpha^{-1/2}\sqrt{\exp(2\lambda T)}(\HN{u_0}+\alpha^{-1/2}\norm{f}_{L^2(I,V^*)})$.

Moreover $\norm{u'}_{L^2(I,V^*)}\leq \norm{Au}_{L^2(I,V^*)}+\norm{f}_{L^2(I,V^*)}\leq \norm{A}\norm{u}_{L^2(I,V)}+\norm{f}_{L^2(I,V^*)}$.

All in all, we obtain:

$$\norm{u}^2_{C([0;T],H)}+\alpha\norm{u}_{L^2(I,V)}^2\leq \exp(2\lambda T)(\HN{u_0}^2+\alpha^{-1}\norm{f}^2_{L^2(I,V^*)})$$

and:

$$\norm{u'}_{L^2(I,V^*)}\leq \norm{A}_{L(V,V^*)}	\alpha^{-1/2}\sqrt{\exp(2\lambda T)}(\HN{u_0}+\alpha^{-1/2}\norm{f}_{L^2(I,V^*)})+\norm{f}_{L^2(I,V^*)}$$

\underline{More regularity}

We tie back to page 25 of \cite{gilardi}. In particular:

\begin{align*}
\int_0^t\HN{u_n'}^2+\int_0^t\langle A u_n, u_n'\rangle_{V^*,V}=\int_0^t(f_1,u_n')_H+\int_0^t \langle f_2, u_n'\rangle_{V^*,V}
\end{align*}

%Now, taking $u_{n0}$ to be the orthogonal projection of $u_0$ onto $V_n$ with the scalar product of $V$, we can assume $\VN{u_{n0}}\leq\VN{u_0}$.

Then:
\begin{align*}
\int_0^t\langle A u_n, u_n'\rangle_{V^*,V}\geq \frac{\alpha}{2}\VN{u_n(t)}^2-\frac{\lambda}{2}\HN{u_n(t)}^2-\frac{\norm{A}}{2}\VN{u_{n0}}
\end{align*}

whereas, as in the proof of \cref{thm:reg_time}:
\begin{align*}
	\left | \int_0^t \langle f_2,u_n' \rangle_{V^*,V}\right | \leq
	\frac{1}{2}\norm{f_2'}_{L^2(I,V^*)}^2 + \frac{1}{2}\norm{u_n}_{L^2(I,V)}^2 + \frac{\alpha}{4}\VN{u_n(t)}^2 +\\
	+ \frac{4}{\alpha}\norm{f_2}_{L^\infty(I,V^*)}^2+ \frac{1}{2}\norm{f_2}_{L^\infty(I,V^*)}^2+\frac{1}{2}\VN{u_{n0}}^2
\end{align*}

Also:
\begin{align*}
\int_0^t(f_1,u_n')_H\leq \frac{1}{2}\norm{f_1}_{L^2(I,H)}^2+\frac{1}{2}\int_0^t\HN{u_n'}^2
\end{align*}

Putting all together:

\begin{align*}
\int_0^t\HN{u_n'}^2+\frac{\alpha}{2}\VN{u_n(t)}^2-\frac{\lambda}{2}\HN{u_n(t)}^2-\frac{\norm{A}}{2}\VN{u_{n0}} \\
\frac{1}{2}\norm{f_2'}_{L^2(I,V^*)}^2 + \frac{1}{2}\norm{u_n}_{L^2(I,V)}^2 + \frac{\alpha}{4}\VN{u_n(t)}^2 +\\
+ \frac{4}{\alpha}\norm{f_2}_{L^\infty(I,V^*)}^2+ \frac{1}{2}\norm{f_2}_{L^\infty(I,V^*)}^2+\frac{1}{2}\VN{u_{n0}}^2+\\
+ \frac{1}{2}\norm{f_1}_{L^2(I,H)}^2+\frac{1}{2}\int_0^t\HN{u_n'}^2
\end{align*}

which brings us to:

\begin{align}
\label{eqn:weak_der_bound}
\frac{1}{2}\int_0^t\HN{u_n'}^2+\frac{\alpha}{4}\VN{u_n(t)}^2-\frac{\lambda}{2}\HN{u_n(t)}^2\leq \\
\frac{1}{2}\norm{f_2'}_{L^2(I,V^*)}^2 + \frac{1}{2}\norm{u_n}_{L^2(I,V)}^2 +\\
+ \frac{8+\alpha}{2\alpha}\norm{f_2}_{L^\infty(I,V^*)}^2+\frac{1+\norm{A}}{2}\VN{u_{n0}}^2+\\
+\frac{1}{2}\norm{f_1}_{L^2(I,H)}^2
\end{align}

and thus, because norms are lower semicontinuous and because we have weak convergence of the time derivative, and $V$-strong convergence of the initial data:

\begin{align*}
\frac{1}{2}\int_0^T\HN{u'}^2\leq 
\frac{1}{2}\norm{f_2'}_{L^2(I,V^*)}^2 
+ \frac{8+\alpha}{2\alpha}\norm{f_2}_{L^\infty(I,V^*)}^2+\frac{1+\norm{A}}{2}\VN{u_{0}}^2+
\frac{1}{2}\norm{f_1}_{L^2(I,H)}^2+\\
\text{limsup}_n \left ( \frac{\lambda}{2}\norm{u_n}_{C([0,T],H)}^2 + \frac{1}{2}\norm{u_n}_{L^2(I,V)}^2 \right )
\end{align*}

Using a purely numeric constant $C$  without dependences on the problem we can write:

\begin{align*}
\int_0^T\HN{u'}^2\leq 
\norm{f_2'}_{L^2(I,V^*)}^2+C(1+\alpha^{-1})\norm{f_2}_{L^\infty(I,V^*)}^2+C(1+\norm{A})\VN{u_{0}}^2+\norm{f_1}_{L^2(I,H)}^2+\\
C\text{limsup}_n \left ( \frac{\lambda}{2}\norm{u_n}_{C([0,T],H)}^2 + \frac{1}{2}\norm{u_n}_{L^2(I,V)}^2 \right )
\end{align*}


For the last term, employing the exact argument as in the first part of the proof:

\begin{align}
\label{eqn:limsup}
\text{limsup}_n \left ( \frac{\lambda}{2}\norm{u_n}_{C([0,T],H)}^2 + \frac{1}{2}\norm{u_n}_{L^2(I,V)}^2 \right )\leq\\
2^{-1}\max(1,\lambda)\max(1,\alpha^{-1}) \text{limsup}_n \left ( \norm{u_n}_{C([0,T],H)}^2 + {\alpha}\norm{u_n}_{L^2(I,V)}^2 \right )\leq\\
2^{-1}\max(1,\lambda)\max(1,\alpha^{-1})\exp(2\lambda T)(\HN{u_0}^2+\alpha^{-1}\norm{f_1+f_2}^2_{L^2(I,V^*)}) \leq \\
2^{-1}\max(1,\lambda)\max(1,\alpha^{-1})\exp(2\lambda T)(\HN{u_0}^2+2\alpha^{-1}\norm{f_1}^2_{L^2(I,V^*)}+2\alpha^{-1}\norm{f_2}^2_{L^2(I,V^*)}) \leq \\
C C_0(\HN{u_0}^2+\alpha^{-1}\norm{f_1}^2_{L^2(I,V^*)}+\alpha^{-1}\norm{f_2}^2_{L^2(I,V^*)})
\end{align}


where $C_0 = \ds 2^{-1}\max(1,\lambda)\max(1,\alpha^{-1})\exp(2\lambda T)$ and $C$ is a purely numeric constant without dependences on the problem.

Therefore:

\begin{align*}
C\int_0^T\HN{u'}^2\leq 
\norm{f_2'}_{L^2(I,V^*)}^2+(1+\alpha^{-1})\norm{f_2}_{L^\infty(I,V^*)}^2+(1+\norm{A})\VN{u_{0}}^2+\norm{f_1}_{L^2(I,H)}^2+\\
C_0(\HN{u_0}^2+\alpha^{-1}\norm{f_1}^2_{L^2(I,V^*)}+\alpha^{-1}\norm{f_2}^2_{L^2(I,V^*)})
\end{align*}


The embedding $H^1(I,V^*)\emb C([0,T],V^*)$ has norm that only depends on $T$, which follows from the equality $f_2(t)=f_2(s)+\int_s^tf_2'$, for $0\leq s \leq t \leq T$, a bound being $1+T$.

Thus:

\begin{align*}
C\int_0^T\HN{u'}^2\leq 
(1+(1+C_0)\alpha^{-1})\norm{f_2}_{H^1(I,V^*)}^2+\\
(1+\norm{A})\VN{u_{0}}^2+C_0\HN{u_0}^2+\\
\norm{f_1}_{L^2(I,H)}^2+C_0\alpha^{-1}\norm{f_1}^2_{L^2(I,V^*)}
\end{align*}


\end{mproof}

\begin{ass}[Even more assumptions for even more regularity]
\label{ass:zero_comp}
We make the hypothesis $u_0=0$, which is the only case that is of interest for us, together with $A$ symmetric. Moreover, we require that the source term is a generic $f\in H^1(I,V^*)$ (in particular, the split $f=f_1 +f_2 \in L^2(I,H)+H^1(I,V^*)$ is not sufficient anymore), so that it is continuous  (\cite{evans}, page 286, theorem 2), and we can therefore ask the additional condition $f(0) \in H$.
\end{ass}

\begin{prop}[More smoothness from zero-order compatibility]
Let \cref{ass:basic_par} and \cref{ass:zero_comp} hold. Then $u \in H^1(I,V)$ (in particular, $u \in C([0,T];V)$), and  $u' \in L^\infty(I,H)$ with the estimates:

\begin{align*}
\norm{u'}^2_{L^\infty(I,H)} \leq \exp(2\lambda T) \left (  \HN{f(0)}^2 +\frac{1}{\alpha} \int_0^T \VSN{ f'}^2\right )
\end{align*}

and:

\begin{align*}
\int_0^T \VN{u'}^2 \leq  \frac{1}{\alpha } \HN{f(0)}^2 + C_1\norm{f}_{H^1(I,V^*)}^2  + 2\lambda\frac{8+\alpha}{\alpha^2}\norm{f}_{L^\infty(I,V^*)}^2
\end{align*}

with $C_0 = \ds 2^{-1}\max(1,\lambda)\max(1,\alpha^{-1})\exp(2\lambda T)$, $C$ is a purely numeric constant without dependences on the problem, $C_1:=\alpha^{-2} + 4 \lambda C C_0 \alpha^{-2} + 2\lambda\alpha ^{-1} $.


\end{prop}

\begin{mproof}

We consider again the functions $u_n$ as defined in \cite{gilardi}, page 22, equation (54).

Thanks to the smoothness of $u_n$ and $f$, which are both $H^1$ in time, we can conclude that $u_n \in H^2(I,V_n)$, due to a bootstrapping argument from equation (56) of \cite{gilardi}.

So, upon differentiation, we get that $\langle u_n'', v_n\rangle_{V^*,V} + \langle A u_n', v_n\rangle_{V^*,V} = \langle f', v_n \rangle_{V^*,V}$ for almost every $t$ (actually, for all $t$). Anyhow, thanks to the smoothness of $u_n'$, we can substitute it as $v_n \in V_n$ and integrate from $0$ to $t$ to obtain:

\begin{align*}
\int_0^t \langle u_n'', u_n'\rangle_{V^*,V} + \int_0^t\langle A u_n',u_n'\rangle_{V^*,V} = \int_0^t \langle f', u_n' \rangle_{V^*,V}
\end{align*}

Equivalently, being $u_n'(t) \in H$, we can write:

\begin{align*}
\int_0^t (u_n'', u_n')_H + \int_0^t\langle A u_n',u_n'\rangle_{V^*,V} = \int_0^t \langle f', u_n' \rangle_{V^*,V}
\end{align*}

Using the assumptions on $A$:

\begin{align*}
\int_0^t (u_n'', u_n')_H + \alpha \int_0^t \VN{u_n'}^2 - \lambda \int_0^t \HN{u_n'}^2 \leq \int_0^t \langle f', u_n' \rangle_{V^*,V}
\end{align*}

Therefore:

\begin{align*}
\frac{1}{2} \HN{u_n'(t)}^2 + \alpha \int_0^t \VN{u_n'}^2 \leq \frac{1}{2} \HN{u_n'(0)}^2+ \lambda \int_0^t \HN{u_n'}^2 +\frac{1}{2\alpha} \int_0^t \VSN{ f'}^2+ \frac{\alpha}{2}\int_0^t \VN{u_n'}^2
\end{align*}

We need to estimate the $H$ norm of $u_n'(0)$, for which the compatibility condition on $f(0)$ is essential.

In fact, because the ODE for $u_n$ held a.e., and now that $u_n' \in H^1(I,V)$, the ODE holds for all times and we can conclude that $(u_n'(0),u_n'(0))_H + \langle A u_n(0), u_n'(0)\rangle_{V^*,V} = ( f(0), u_n'(0))_H$. Here comes in handy the fact that $u_0=0$, so that $u_n(0)=0$ too and therefore $\HN{u_n'(0)}^2  \leq 2^{-1}\HN{f(0)}^2 + 2^{-1} \HN{u_n'(0)}^2$ and therefore, $\HN{u_n'(0)}^2\leq \HN{f(0)}$.

With this bound:

\begin{align*}
\HN{u_n'(t)}^2 + \alpha \int_0^t \VN{u_n'}^2 \leq \HN{f(0)}^2+ 2\lambda \int_0^t \HN{u_n'}^2 +\frac{1}{\alpha} \int_0^T \VSN{ f'}^2
\end{align*}

\underline{The $L^\infty$ estimate}

Gronwall's lemma (in the form of \cite{gilardi} at page 19) yields:

\begin{align*}
\HN{u_n'(t)}^2 \leq \exp(2\lambda T) \left (  \HN{f(0)}^2 +\frac{1}{\alpha} \int_0^T \VSN{ f'}^2\right )
\end{align*}

This alone doesn't show that $u'\in L^\infty(I,H)$, but if this was true, and is $u_n'(t)\rightarrow_H u'(t)$ for a.e. $t$ modulo subsequences, then this would yield the required estimate on the norm.

Now, $u_n, u_n' \rightharpoonup_H u, u'$ (where $u'$ is the $L^2(I,H)$ representative of the derivative of $u$ in the distributional sense, see the proof of \cref{thm:reg_time}). This is true modulo a common subsequence

%, but since we also know the boundedness of $u_n, u_n' $ in $L^2(I,H)$ (cfr the proof of \cref{thm:reg_time} and page 23 of \cite{gilardi}), we conclude the convergence of the full sequences.

Because $u_n', u'$ are the $L^2(I,H)$ sense derivatives of $u_n, u$ (thanks to the injectivity of $H^* \emb V^*$), then, by the compactness theorem 2.1 at page 271 of \cite{navier_stokes} with $X_0=X=X_1=H$ we conclude that $u_n \rightarrow u$ strongly in $L^2(I,H)$. Thanks to proposition 2.13 at page 10 of \cite{kreuter}, $u_n\rightarrow_H u$ for a.e. $t$, modulo a further subsequence.

The bound shown above implies that $u' \in L^\infty(I,H)$ as we wished. 

\underline{$u' \in L^2(I,V)$}


We know:


\begin{align*}
\HN{u_n'(t)}^2 + \alpha \int_0^t \VN{u_n'}^2 \leq \frac{1}{2} \HN{f(0)}^2+ 2\lambda \int_0^t \HN{u_n'}^2 +\frac{1}{\alpha} \int_0^T \VSN{ f'}^2
\end{align*}

so that:


\begin{align*}
\alpha \int_0^T \VN{u_n'}^2 \leq \frac{1}{2} \HN{f(0)}^2+ 2\lambda \int_0^T \HN{u_n'}^2 +\frac{1}{\alpha} \int_0^T \VSN{ f'}^2
\end{align*}

But the proof of \cref{thm:const_track} says (cfr. \cref{eqn:weak_der_bound}):

\begin{align*}
\frac{1}{2}\int_0^t\HN{u_n'}^2+\frac{\alpha}{4}\VN{u_n(t)}^2\leq \\
\frac{\lambda}{2}\HN{u_n(t)}^2 + \frac{1}{2}\norm{f'}_{L^2(I,V^*)}^2 + \frac{1}{2}\norm{u_n}_{L^2(I,V)}^2 + \frac{8+\alpha}{2\alpha}\norm{f}_{L^\infty(I,V^*)}^2\\
\end{align*}

In particular:

\begin{align*}
2\lambda\int_0^T\HN{u_n'}^2\leq \\
4 \lambda \left (\frac{\lambda}{2} \norm{u_n}^2_{C([0,T];H)} + \frac{1}{2}\norm{u_n}_{L^2(I,V)}^2 \right )+ 2\lambda \norm{f'}_{L^2(I,V^*)}^2  + 2\lambda\frac{8+\alpha}{\alpha}\norm{f}_{L^\infty(I,V^*)}^2\\
\end{align*}

so that:

\begin{align*}
\alpha \int_0^T \VN{u_n'}^2 \leq \frac{1}{2} \HN{f(0)}^2 +\frac{1}{\alpha} \int_0^T \VSN{ f'}^2+\\
4 \lambda \left (\frac{\lambda}{2} \norm{u_n}^2_{C([0,T];H)} + \frac{1}{2}\norm{u_n}_{L^2(I,V)}^2 \right )+ 2\lambda \norm{f'}_{L^2(I,V^*)}^2  + 2\lambda\frac{8+\alpha}{\alpha}\norm{f}_{L^\infty(I,V^*)}^2
\end{align*}

As we have already seen (or as it is explained at page 23 of \cite{gilardi}), the term in the brackets is bounded above.

%Now, remember that $u_n$ was bounded in $L^2(I,V)$, and that any weakly convergent subsequence would solve the parabolic problem. Thus, $u_n \rightharpoonup u$ in $L^2(I,V)$ (cfr. page 23 in \cite{gilardi}).

Morevoer, we showed in the proof of \cref{thm:reg_time} that $u_n'$ is bounded in $L^2(I,H)$, and that any weakly convergent subsequence would reach $u'$, thus yielding the weak convergence of the full sequence.


Also we now have that $u_n'$ is bounded in $L^2(I,V)$, thus admitting a weakly convergent subsequence to $w$, in $L^2(I,V)$. Because it is implied the weak convergence in $L^2(I,H)$, to $w$, we obtain that $w=u' $. Reiterating the argument and thanks to the boundedness of the full sequence, the full sequence must converge weakly in $L^2(I,V)$ to $u'$.

By lower semicontinuity of the norm:

\begin{align*}
\alpha \int_0^T \VN{u'}^2 \leq \HN{f(0)}^2 +\frac{1}{\alpha} \int_0^T \VSN{ f'}^2+\\
4 \lambda \limsup_n\left (\frac{\lambda}{2} \norm{u_n}^2_{C([0,T];H)} + \frac{1}{2}\norm{u_n}_{L^2(I,V)}^2 \right )+ 2\lambda \norm{f'}_{L^2(I,V^*)}^2  + 2\lambda\frac{8+\alpha}{\alpha}\norm{f}_{L^\infty(I,V^*)}^2
\end{align*}

Remembering \cref{eqn:limsup}:

\begin{align*}
\text{limsup}_n \left ( \frac{\lambda}{2}\norm{u_n}_{C([0,T],H)}^2 + \frac{1}{2}\norm{u_n}_{L^2(I,V)}^2 \right )\leq\\
C C_0 \alpha^{-1}\norm{f}^2_{L^2(I,V^*)}
\end{align*}


where $C_0 = \ds 2^{-1}\max(1,\lambda)\max(1,\alpha^{-1})\exp(2\lambda T)$ and $C$ is a purely numeric constant without dependences on the problem, so that:

\begin{align*}
\int_0^T \VN{u'}^2 \leq \\ \frac{1}{\alpha } \HN{f(0)}^2 + C_1\norm{f}_{H^1(I,V^*)}^2  + 2\lambda\frac{8+\alpha}{\alpha^2}\norm{f}_{L^\infty(I,V^*)}^2
\end{align*}

for $C_1:=\alpha^{-2} + 4 \lambda C C_0 \alpha^{-2} + 2\lambda\alpha ^{-1} $.

Because the embedding $H\emb V$ is injective, $u'$ is also the weak derivative, $L^2(I,V)$ sense, of $u \in L^2(I,V)$, so that $u \in H^1(I,V)$. And thus, thanks to theorem 2 at page 286 of \cite{evans}, we obtain $u \in C([0,T]; V)$.

\end{mproof}


\section{Application to inhomogeneous parabolic problems}

\subsection{Inhomogeneous Dirichlet problem}
\label{subs:inh_diri}

We make the following assumption.

\begin{ass}[Assumptions for \cref{pb:diri}]
\label{ass:diri}
We assume $\Omega \cc D $ to be bounded Lipschitz domains, so that $U:=D\setminus \Omega$ is bounded Lipschitz too and  the trace operator is bounded surjective onto $H^{1/2}(\partial U)$, with a right inverse $E$ (see theorem 3.37 at page 102 of \cite{mclean}). For such a choice we also have $H^1_0=H^1\cap \text{ker }\tr$, see \cite{leoni}, page 595, theorem 18.7.

Moreover, we select $f \in H^1(I, H^{1/2}(\Gamma_f))$, $f(0)=0$.
\end{ass}

Note that, given a bounded extension operator $E: H^{1/2}(\partial U) \rightarrow H^1(U)$, we obtain by \cref{lemma:bochner_Hk_map} that $Ef \in H^2(I, H^1(U))$. We have defined $\tr u (t):= tr(u(t))$ and analogously $Eu(t):=E(u(t))$ (see \cref{prop:trace}).

Call $H=L^2(U)$, $V=\{ v \in H^1(U), \tr u = 0 \text{ on } \Gamma_m\}=:H^1_c$. $V$ is a closed subspace of $H^1$, which is Hilbert separable, hence also Hilbert separable. We norm it with the full $H^1$ norm. Because $H^1_0(U)$ is dense in $H$, so is $V$ and we obtain a Gelfand triple. That $V$ is a closed subspace of $H^1$ follows from the observation that if $u_n\rightarrow u$ in the $V$ norm, then $\tr u_n \rightarrow \tr u$ in $L^2(\partial U)$. We can take an almost everywhere pointwise convergent sequence, so that $\tr u_n \rightarrow \tr u$ a.e., and by the fact that $\Gamma_m$ has positive Hausdorff measure, we conclude $\tr u = 0$ on $\Gamma_m$.

We define $A:= H^1 \rightarrow H^{1*}$ by $(Au)v:=\int_u\nabla u \nabla v$. This operator can be the recast to $V\rightarrow H^{-1}$ and $V\rightarrow V^*$.

The problem under consideration is the following. For $U = D\setminus \Omega$ we have:

\begin{pb}[Inhomogeneous heat equation, Dirichlet conditions]
\label{pb:diri}
\begin{align}
u_t - \Delta u = 0 \text{ in } (0,T)\times U\\
u(\Sigma_f)=f\\
u(\Sigma_m)=0\\
u(0)=0
\end{align}

By this we mean:

\begin{align}
u \in W(I,H^1_c) \\
u_t|_{H^{-1}} + A u = 0 \text{ in }H^{-1} \text{ and for a.e. } t \in (0,T) \\
\tr u = f \text{ on } \Sigma_f\\
u(0)=0
\end{align}

\end{pb}

\begin{thm}[Well posedness and regularity for \cref{pb:diri}]
\label{prop:diri_wp}
Given \cref{ass:diri}, the solution $u$ to \cref{pb:diri} is unique with $u_t \in L^2(I,H)$. 

%The uniqueness holds in more generally in $L^2(I,V)\cap H^1(I, H^{-1})$.

The problem is equivalent to:

\begin{pb}[Equivalent formulation with extension]
\label{pb:diri_ext}
\begin{align}
u_0 \in W(I,H^1_0) \\
u_0' + A u_0 = -((\bar{u}',\cdot)_H+A \bar{u}) \text{ in }H^{-1} \text{ and for a.e. } t \in (0,T) \\
u_0(0)=0
\end{align}
\end{pb}

with $\bar{u}$ any given $\bar{u}\in H^1(I,H^1_c(U))$ such that $\tr \bar{u} =f$ on $\Sigma_f$, and with $\bar{u}(0)=0$. This means that $u$ solves \cref{pb:diri} $\implies$ $u-\bar{u}$ solves \cref{pb:diri_ext}, and if $u_0(\bar{u})$ solves  \cref{pb:diri_ext}, then $\bar{u}+u_0(\bar{u})$ solves \cref{pb:diri}.

Furthermore: 

\begin{align}
\norm{u}^2_{C([0;T],H)}+\norm{u}_{L^2(I,H)}^2+ \norm{\nabla u}_{L^2(I,H)}^2 + \norm{u'}^2_{L^2(I,H)}\leq C(T)\norm{\bar{u}}_{H^1(I,V)}^2
\end{align}

with $C>1$, only dependent on $T$, smoothly, exploding for large $T$.

\end{thm}
\begin{mproof}

\underline{Extension of the boundary data}

Let $\bar{u}\in H^1(I,H^1_c(U))$ be such that $\tr \bar{u} =f$ on $\Sigma_f$, and with $\bar{u}(0)=0$. We can choose for instance $E\tilde{f}$, see \cref{prop:trace}, where $\tilde{f}=0$ on $\Sigma_m$, $\tilde{f}=f$ on $\Sigma_f$. $\tilde{f} \in H^1(I,H^{1/2}(\partial U))$, because $\Gamma_f$ and $\Gamma_m$ have positive distance (see the definition of the norm in \cite{grisvard}, page 20).  

\underline{Reformulation (first part)}

Consider the following commutative diagram, where $V = H^1_c$, $W=H^1_0$:

\[\begin{tikzcd}
	{H^1_c} &&& {(H^1_c)^*} \\
	& H & {H^*} \\
	{H^1_0} &&& {H^{-1}}
	\arrow["c", hook, from=3-1, to=1-1]
	\arrow["a", hook, from=1-1, to=2-2]
	\arrow["b"', hook, from=3-1, to=2-2]
	\arrow["{a^*}", hook, from=2-3, to=1-4]
	\arrow["r", from=2-2, to=2-3]
	\arrow["{b^*}"', hook, from=2-3, to=3-4]
	\arrow["{c^*}", hook, from=1-4, to=3-4]
\end{tikzcd}\]

Here, $a,b,c$ are the trivial injections, $r$ the Riesz isomorphism $h\mapsto(h,\cdot)_H$. 

Now $(i_W (u-\bar{u}) )'+A(u - \bar{u}) = (i_Vu)'|_{H^{-1}}-(i_V\bar{u})'|_{H^{-1}}+Au - A\bar{u} = \ind{\cref{prop:sanity}}  = (i_Vu)'|_{H^{-1}}-(i_V\bar{u}_t)|_{H^{-1}}+Au - A\bar{u} = -(i_V\bar{u}_t)|_{H^{-1}}- A\bar{u}$ if $u$ solves \cref{pb:diri}, where $\bar{u}_t$ is the weak derivative of $\bar{u}$ in the $H^1(I,V)$ sense. Call $u_0 = u-\bar{u}$. By again \cref{prop:sanity}, $u_0 \in W(I,H^1_0)$.

%That $\tr E \tilde{f}=\tilde{f}$ is easily seen for a.a. $t$, and by \cref{lemma:bochner_Hk_map}, we get $\bar{u} \in H^1(I,H^1_c(U))$. The last result could be applied since $H^{1/2}(\partial U)$ is the continuous image of a separable space, hence it is separable, and it is also Banach (again by \cite{grisvard}, page 20).
%
%
%The initial value of $\bar{u}$ is null. To see this, we know that the unique time-continuous representative $U\tilde{f}$ of $\tilde{f}$ (\cite{evans}, page 286, theorem 2, and \cref{prop:cts_repr}) satisfies $\lim_{t\rightarrow 0^+}U\tilde{f}(t)=0$, the limit being in $H^{1/2}(\partial U)$. We also know that $\bar{u}$ has a unique continuous representative $U\bar{u}$ by similar arguments, and $U\bar{u}(t) \rightarrow U\bar{u}(0)$ for $t\rightarrow 0^+$. Now, $ [U\tilde{f}] = \tilde{f} = \tr \bar{u} =[ \tr U\bar{u} ]$, where the square brackets are equivalence classes in time. Thus $U\tilde{f} =  \tr U\bar{u}$ for all $t$, and thus $\tr U\bar{u}\rightarrow 0$ for small times. Applying $E$ we get $U\bar{u}\rightarrow 0$ which is what we wanted to show.

This motivates us to consider the problem: 

\begin{align}
u_0 \in W(I,H^1_0) \\
u_0' + A u_0 = -(f_1+f_2) \text{ in }H^{-1} \text{ and for a.e. } t \in (0,T) \\
u_0(0)=0
\end{align}

Here, $f_1:=(i_V\bar{u}_t)|_{H^{-1}}=c^*a^*ra\bar{u}_t = b^*r (a\bar{u}_t) \in L^2(I,H)$.

Moreover, $A \in L(V, H^{-1})$, so, by \cref{lemma:bochner_Hk_map}, $f_2:=A\bar{u} \in H^1(I, H^{-1})$.

\underline{Existence}

By \cref{thm:par_reg} we get a solution of the above problem with $u_0' \in L^2(I, H)$.

And now, let $u:=\bar{u}+c u_0 =\bar{u}+ u_0 $. We claim it is a solution. The initial and boundary conditions are surely satisfied. We check it is in $W(I,V)$ and is satisfies the partial differential equation.

By \cref{prop:sanity}, we have both $\bar{u}, c u_0 \in W(I,V)$. The derivative of $\bar{u}$ becomes $i_V \bar{u}_t$, see \cref{prop:sanity}. Therefore $(i_V(\bar{u}+c u_0))'|_{H^1_0} = c^* (i_V(\bar{u}+c u_0))' = c^*i_V \bar{u}_t + c^*i_V(c u_0)' = b^*r (a\bar{u}_t) + i_W(u_0)'$ by \cref{prop:sanity}.

Using the pde of $u_0$, $... = b^*r (a\bar{u}_t) - Au_0 -f_1 -f_2 = -A(u_0+\bar{u})$.

\underline{Uniqueness}

For two solutions $u_1, u_2$ of $\ref{pb:diri}$ we can form $d:=u_1-u_2\in W(I,H^1_0)$ by \cref{prop:sanity}. Clearly, $d(0)=0$. Moreover, $(i_{H^1_0} d)' = \ind{\cref{prop:sanity}}=(i_V u_1)'|_{H^1_0}-(i_V u_2)'|_{H^1_0} = A(u_1-u_2)$.

By uniqueness stated in \cref{thm:well_pos_parabolic} we obtain $d=0$ in $L^2(I,H)$, so that the solution is unique and doesn't depend on the choice of the extension of the Dirichlet datum.

\underline{Reformulation (part 2)}

Therefore $u=\bar{u}+u_0$ above is the unique solution of \cref{pb:diri}. So, given any $\bar{u}\in H^1(I,H^1_c(U))$ such that $\tr \bar{u} =f$ on $\Sigma_f$, and with $\bar{u}(0)=0$, we can construct $u_0$ as above and get $u=\bar{u}+u_0$ solving \cref{pb:diri}.

Viceversa, let $u$ solve \cref{pb:diri}. Call $u_0 = u- \bar{u}$. Then, as seen above, $u_0 \in W(I,H^1_0)$ and $(i_V (u-\bar{u}) )'|_{H^{-1}}+A(u - \bar{u}) = (i_Vu)'|_{H^{-1}}-(i_V\bar{u})'|_{H^{-1}}+Au - A\bar{u} = \ind{\cref{prop:sanity}}  = (i_Vu)'|_{H^{-1}}-(i_V\bar{u}_t)|_{H^{-1}}+Au - A\bar{u} = -(i_V\bar{u}_t)|_{H^{-1}}- A\bar{u}$ if $u$ solves \cref{pb:diri}, where $\bar{u}_t$ is the weak derivative of $\bar{u}$ in the $H^1(I,V)$ sense. Call $u_0 = u-\bar{u}$. By again \cref{prop:sanity}, $(i_W (u-\bar{u}) )'+A(u - \bar{u}) = (i_Vu)'|_{H^{-1}}-(i_V\bar{u})'|_{H^{-1}}+Au - A\bar{u} = (i_Vu)'|_{H^{-1}}-(i_V\bar{u}_t)|_{H^{-1}}+Au - A\bar{u} = -(i_V\bar{u}_t)|_{H^{-1}}- A\bar{u} = -b^*r(a\bar{u}_t)-A\bar{u}$. Moreover, $u_0(0)=0$, so that $u_0$ solves \cref{pb:diri_ext}.

%\begin{itemize}
%\item $u_0 \in L^2(I,H^1_0)$
%\item $u_0' \in L^2(I,V^*)+L^2(I,H)\subseteq L^2(I,H^{-1})$
%\item $u_0(0)=0$ in $H$
%\item $u_0$ solves $u_0' + A u_0 = -((\bar{u}',\cdot)_H+\langle A \bar{u},\cdot\rangle_{H^{-1},H^1_0})$
%\end{itemize}

\underline{Regularity}

Let $u=\bar{u}+u_0$ be the unique solution, as before, of \cref{pb:diri}. From \cref{prop:sanity} we know $(i_V(\bar{u}))'=i_V(\bar{u}_t)=a^*r(a\bar{u}_t)$, and $i_V(cu_0)'=a^*r (u_0')$, for $u_0' \in L^2(I,H)$ the representative of $(i_W(u_0))'$, equivalently, the weak derivative of $u_0$ in the $H^1(I,H)$ sense. It follows that $(i_Vu)' = a^*r(a\bar{u}_t+u_0')$, proving the additional time smoothness claim.

\underline{Stability}

Let $\bar{u}\in H^1(I,H^1_c(U))$ such that $\tr \bar{u} =f$ on $\Sigma_f$, and with $\bar{u}(0)=0$. Consider $u_0$. Then, by \ref{thm:const_track}:


\begin{align*}
\norm{u_0}^2_{C([0;T],H)}+\alpha\norm{u_0}_{L^2(I,H^1_0)}^2\leq \exp(2\lambda T)\alpha^{-1}\norm{(\bar{u}',\cdot)_H+ A \bar{u}}^2_{L^2(I,H^{-1})}
\end{align*}

\begin{align*}
C\norm{u_0'}^2_{L^2(I,H)}\leq 
(1+(1+C_0)\alpha^{-1})\norm{A \bar{u}}^2_{H^1(I, H^{-1})}+\\
\norm{(\bar{u}',\cdot)_H}_{L^2(I,H)}^2+C_0\alpha^{-1}\norm{(\bar{u}',\cdot)_H}^2_{L^2(I,H^{-1})}
\end{align*}

$C_0 = \ds 2^{-1}\max(1,\lambda)\max(1,\alpha^{-1})\exp(2\lambda T)$.

We norm $H^1_0$ with the full $H^1$ norm too. Then:
\begin{align*}
\norm{(\bar{u}',\cdot)_H+ A \bar{u}}_{L^2(I,H^{-1})}\leq \\
\sup_{\norm{v}_{L^2(I,H^{1}_0)}=1}\norm{\bar{u}'}_{L^2(I,H)}\norm{v}_{L^2(I,H)}+\norm{\nabla \bar{u}}_{L^2(I,H)}\norm{\nabla v}_{L^2(I,H)}\leq \\C(\norm{\bar{u}'}_{L^2(I,H)}+\norm{\nabla \bar{u}}_{L^2(I,H)})
\end{align*}

By \cref{lemma:bochner_Hk_map}, $\norm{A \bar{u}}_{H^1(I, H^{-1})}\leq \norm{A}_{L(V,H^{-1})}\norm{\bar{u}}_{H^1(I,V)}$ (we could apply it since $H^{-1}$ is separable, as a dual of a reflexive Banach space).

Finally, $\norm{(\bar{u}',\cdot)_H}^2_{L^2(I,H^{-1})}\leq \norm{\bar{u}'}_{L^2(I,H)}^2$.

We can then say:

\begin{align*}
\norm{u_0}^2_{C([0;T],H)}+C\alpha\norm{u_0}_{L^2(I,H^1_0)}^2\leq \exp(2\lambda T)\alpha^{-1}\norm{\bar{u}}^2_{H^1(I,V)}\\
C\norm{u_0'}^2_{L^2(I,H)}\leq 
((1+(1+C_0)\alpha^{-1})\norm{A}_{L(V,H^{-1})}^2 +1 +C_0\alpha^{-1})\norm{\bar{u}}^2_{H^1(I,V)}
\end{align*}

Now, $\langle Av, v\rangle_{H^{-1}, H^1_0} + 1\cdot  \HN{v}^2 = 1\cdot \norm{v}_{H^1_0}^2$, so that $\alpha=\lambda=1$. Moreover, $\langle Au, v\rangle_{H^{-1}, H^1_0}\leq \VN{u}\norm{v}_{H^1_0}$, i.e. $\norm{A}_{L(V,H^{-1})}\leq 1$.

Therefore $\norm{u_0}^2_{C([0;T],H)}+\norm{u_0}_{L^2(I,H^1_0)}^2 + \norm{u_0'}^2_{L^2(I,H)}\leq C(T)\norm{\bar{u}}_{H^1(I,V)}^2$ with $C>1$, only dependent on $T$, smoothly, exploding for large $T$.

Now, let's analyse the norms of $\bar{u}$. Because $\bar{u}\in H^1(I,V)$, then, $\bar{u} \in C([0,T],V)\emb C([0,T],H)$, where the embedding is non-expansive by the choice of the norm of $V$. Therefore $\norm{\bar{u}}_{C([0;T],H)}\leq \norm{\bar{u}}_{C([0;T],V)}\leq (1+T)\norm{\bar{u}}_{H^1(I,V)}$. We can therefore conclude that $\norm{u}^2_{C([0;T],H)}+\norm{u}_{L^2(I,H^1_0)}^2 + \norm{u'}^2_{L^2(I,H)}\leq C(T)\norm{\bar{u}}_{H^1(I,V)}^2$ with $C>1$, only dependent on $T$, smoothly, exploding for large $T$.



\end{mproof}

\subsection{Inhomogeneous Neumann-Dirichlet problem}

We make the following assumption.

\begin{ass}[Assumptions for \cref{pb:diri}]
\label{ass:neu}
We keep \cref{ass:diri} (apart from the Dirichlet datum). We consired $g \in H^1(I, L^2(\Gamma_f))$, $g(0)=0$.
\end{ass}

Again, call $H=L^2(U)$, $V=\{ v \in H^1(U), \tr u = 0 \text{ on } \Gamma_m\}=:H^1_c$. $H,V$ induce a Gelfand triple as seen before. 


The problem under consideration is:

\begin{pb}[Inhomogeneous heat equation, Neumann conditions]
\label{pb:neu}
\begin{align}
u_t - \Delta u = 0 \text{ in } (0,T)\times U\\
\partial_\nu u(\Sigma_f)=g\\
u(\Sigma_m)=0\\
u(0)=0
\end{align}

By this we mean:

\begin{align}
u \in W(I,H^1_c) \\
u_t + A u = G \text{ in } V^* \text{ and for a.e. } t \in (0,T) \\
u(0)=0
\end{align}

where $\langle G(t), v \rangle_{V^*,V}:=\int_{\Gamma_f} g(t)\tr v d\sigma$, $\sigma$ the $1$-codimensional Hausdorff measure, and $A$ was introduced before in $L(V,H^{-1})$.

\end{pb}

By \cref{lemma:bochner_Hk_map}, $G \in H^1(I,V^*)$. In fact, define $T: L^2(\Sigma_f)\rightarrow V^*$ by $\langle Tg,v \rangle_{V^*,V}:=\int_{\Gamma_f} g\tr v d\sigma$. Then, $\langle Tg,v \rangle_{V^*,V} \leq \norm{g}_{L^2(\Gamma_f)} \VN{v}$ by trace theory. Now, $G(t)=Tg(t)$.

Moreover, $\langle A v, v \rangle_{V^*,V}+ 1 \cdot \HN{v} = 1\cdot \norm{V}$, so that we can immediately conclude:

\begin{thm}[Well posedness and regularity for \cref{pb:neu}]
\label{prop:wp_neu}
Given \cref{ass:neu}, the solution $u$ to \cref{pb:neu} is unique with $u_t \in L^2(I,H)$.

Furthermore: 

\begin{align}
\norm{u}^2_{C([0;T],H)}+\norm{u}_{L^2(I,H)}^2+ \norm{\nabla u}_{L^2(I,H)}^2 + \norm{u'}^2_{L^2(I,H)}\leq C(T)\norm{g}_{H^1(I,L^2(\Gamma_f))}^2
\end{align}

with $C>1$, only dependent on $T$, smoothly, exploding for large $T$.
\end{thm}
\begin{mproof}

It is an application of \cref{thm:well_pos_parabolic}, \cref{thm:reg_time} and \cref{thm:const_track}.
\end{mproof}

\section{Reformulation of parabolic equations}

We just saw that the two parabolic equations of interest can be recasted into the problem of finding $u\in W(I,V)$, $u(0)=0$, $u_t+Au=f$ for a.e. $t$ in $V^*$, with notation from preceding sections.

In particular, $f \in L^2(I, V^*)$ and so is $Au$ (because $A\in L(V,V^*)$, and by \ref{lemma:bochner_Hk_map}).

Call then $E(u):=u_t+Au-f \in L^2(I,V^*)$ and $W_0(I,V)$ the $W(I,V)$ functions with zero initial value. Then, the differential equation reads $\langle E(u)(t),v\rangle_{V^*,V}=0$ for all $v\in V$, for a.a. $t$, equivalently, $E(u)=0$ for a.a. $t$. Thus, we are interested in the abstract problem:

\begin{pb}[Even more abstract parabolic equation]
\label{pb:more_abstr_par}
Given a function $E: W(I,V)\rightarrow L^2(I,V^*)$, find $u\in W_0(I,V)$, such that $E(u)=0$ for a.a. $t$.
\end{pb}
 
We can view $L^2(I,V^*)\cong L^2(I,V)^*$.

Hence $\langle E(u), v\rangle_{L^2(I,V)^*, L^2(I,V)}=\int_I \langle E(u)(t),v(t) \rangle_{V^*,V} dt$ (see \cite{hinze}, theorem 1.31 at page 39).

\color{black}


We are now ready to restric both state and adjoint space.

\begin{defn}[$Q(I,V)$]
We define $Q(I,V)=H^{1,1}=L^2(I,V)\cap H^1(I,H)$, with the norm $\norm{v}_Q^2=\norm{v}_{L^2(I,V)}^2 + \norm{v_t}_{L^2(I,H)}^2$.
\end{defn}

\begin{prop}[Properties of $Q$]
\label{prop:Q}
There holds:
\begin{itemize}
	\item $Q=Q(I,V)$ is Hilbert with $(v,w)_{L^2(I,V)} + (v_t,w_t)_{L^2(I,H)}$ 
	\item $Q(I,V)$ is dense in $L^2(I,V)$
	\item $Q(I,V)\emb C([0,T],H)$
	\item $Q_0(I,V)$ is dense in $L^2(I,V)$, $Q_0(I,V)$ the space of $Q(I,V)$ function with zero initial value
	\item $Q(I,V) = W(I,V)\cap H^1(I,H)$, $Q_0(I,V) =  W_0(I,V)\cap H^1(I,H)$ as sets. There holds that the $W(I,V)$ derivative is represented by the $H^1(I,H)$ derivative and $\langle u', v\rangle_{V^*,V} = (u',v)_H$, with the suitable interpretations of $u'$ (on the left, we have $i_V(u)'$, $i_V$ the Gelfand triple embedding; on the right we have $u$, seen in $H$, and then weakly differentiated in the $H^1(I,H)$ sense)
	\item integration by parts in time holds: $\int_I(v_t,w)_H = -\int_I(w_t,v)_H +(v(T),w(T))_H-(v(0),w(0))_H$
	\item if $q_n$ is bounded in $Q(I,V)$, then there exists a weakly convergent subsequence $q_k$ such that $q_k\weakc q$ in $L^2(I,H)$, $\partial_i q_k\weakc \partial_i q$ in $L^2(I,H)$ and $q_k'\weakc q'$ in $L^2(I,H)$
\end{itemize}
\end{prop}
\begin{mproof}

\underline{Completeness}

We have the inclusions $L^2(I,H)\subseteq L^2(I,V)\cap H^1(I,H)\subseteq H^1(I,V)$.

If $q_n$ is Cauchy in $Q$, then it is Cauchy in the individual norms of $ L^2(I,V), H^1(I,H)$, so that $q_n$ convergens to two limits, one in $L^2(I,V)$ and one in $H^1(I,H)$. The convergence is common in $L^2(I,H)$, which implies that the two limits coincide at $q \in Q(I,V)$. The convergence in $Q(I,V)$ of $q_n$ to $Q$ follows from the individual convergences of $q_n, \nabla q_n, q_{nt}$ in $ L^2(I,H), L^2(I,H),L^2(I,H)$.

\underline{Density}

We have $C_c^\infty(I,V) \subseteq Q(I,V) \subseteq L^2(I,V)$. The first inclusion holds because of \cref{prop:weak_class}, so that $C_c^\infty(I,V)\subseteq H^1(I,V)$. Moreover $H^1(I,V)\subseteq  Q(I,V) $ trivially, where the $H^1(I,H)$ derivative is the $H^1(I,V)$ derivative. $C_c^\infty(I,V)$ is dense in $ L^2(I,V)$ by \cite{hinze}, page 39, lemma 1.9.

\underline{Continuity}

Follows from the embedding $H^1(I,H)\emb C([0,T],H)$, as seen in \cite{evans}, theorem 2 of page 286.

\underline{More density}

We can therefore speak of initial values. In particular,  $C_c^\infty(I,V) \subseteq Q_0(I,V)\subseteq L^2(I,V)$, and as before, the density result follows.

\underline{Relationship with $W(I,V)$}

Consider the chain:

\[\begin{tikzcd}
	V & H & {H^*} & {V^*}
	\arrow["a", hook, from=1-1, to=1-2]
	\arrow["r", from=1-2, to=1-3]
	\arrow["{a^*}", hook, from=1-3, to=1-4]
\end{tikzcd}\]

where $a$ is the trivial embedding and $r$ the Riesz isomorphism.

We claim that for $v \in Q(I,V)$, then $(a^*ra v)' = a^*r (av)'$, where $av \in H^1(I,H)$. In fact, for $\phi \in C^\infty_c(I)$, we get $\int_I a^*ra v \phi' = \ind{\ref{prop:bochner_bound}} = a^*r\int_I av\phi' = -a^*r\int_I(av)'\phi =-\int_I a^*r(av)' \phi$.

Now, let $u \in W(I,V)$, with $(a^*ra u)' = a^*r h$, $h \in L^2(I,H)$. Then $ a^*r \int_I h \phi = \ind{\ref{prop:bochner_bound}} = \int_I a^*r h \phi  = \int_I (a^*ra u)' \phi =  a^*r (-\int au \phi')$.

We know that $a^*$ is injective and so is $r$, so that $ \int_I h \phi  = -\int_I au \phi'$ as we wanted.

\underline{Integration by parts}

We note that for $v,w \in Q(I,V)\subseteq W(I,V)$, we have $\int_I \langle (a^*ra v)', w\rangle_{V^*,V} = \int_I \langle a^*r(a v)', w\rangle_{V^*,V} = \int_I((a v)', aw)_H$. We can now apply theorem 3.11 at page 148 of \cite{trol} to conclude that  $\int_I(v_t,w)_H = -\int_I((aw)_t,av)_H +((av)(T),(aw)(T))_H-((av)(0),(aw)(0))_H$

\underline{Weak convergence}

At first we note that $\partial_i, \partial_t$ are linear bounded operators from $Q(I,V)$ to $L^2(I,H)$.

Remember that in any case, $V$ is a closed subspace of $H^1$. Then, $\partial_i : V\rightarrow H$ is linear and bounded, because $V$ is bounded by the full $H^1$ norm, as we declared already.

Therefore, by \cref{lemma:bochner_Hk_map}, $\partial_i$ extends to a linear bounded map from $L^2(I,V)$ to $L^2(I,H)$, therefore, to a linear bounded map on $Q(I,V)	$, in the sense of:

\[\begin{tikzcd}
	{Q(I,V)} & {L^2(I,V)} & {L^2(I,H)}
	\arrow["i", hook, from=1-1, to=1-2]
	\arrow["{\partial_i}", from=1-2, to=1-3]
\end{tikzcd}\]

Because $q_n$ is bounded in the Hilbert space $Q(I,V)$, it has a weakly convergent subsequence $q_k\weakc q \in Q(I,V)$. Therefore, $\partial_i (i(q_k))\weakc \partial_i (i(q))$ in $L^2(I,H)$. By the Hilbert space property of $L^2(I,H)$ we conclude that $(\partial_i q_k,p)_{L^2(I,H)}\rightarrow (\partial_i q,p)_{L^2(I,H)}$ for all $p \in L^2(I,H)$.

For the time derivative we can draw a similar diagram:

\[\begin{tikzcd}
	{Q(I,V)} & {H^1(I,H)} & {L^2(I,H)}
	\arrow["j", hook, from=1-1, to=1-2]
	\arrow["{\partial_t}", from=1-2, to=1-3]
\end{tikzcd}\]

We therefore obtain $(\partial_t q_k,p)_{L^2(I,H)}\rightarrow (\partial_t q,p)_{L^2(I,H)}$ for all $p \in L^2(I,H)$.

The convergence $(q_k,p)_{L^2(I,H)}\rightarrow (q,p)_{L^2(I,H)}$ for all $p \in L^2(I,H)$ follows analogously.

\end{mproof}

We can therefore restrict the testing space.

\begin{prop}[Equivalent testing]
\label{prop:eq_test}
Let $E: W(I,V)\rightarrow L^2(I,V^*)$, and $u\in W_0(I,V)$.

Then:

\begin{align*}
E(u)=0 \\
\iff \\
\langle E(u), v\rangle_{L^2(I,V)^*, L^2(I,V)}=0 \quad \forall v \in L^2(I,V)\\
\iff \\
\langle E(u), v\rangle_{L^2(I,V)^*, L^2(I,V)}=0 \quad \forall v \in W^0(I,V)\\
\iff \\
\langle E(u), v\rangle_{L^2(I,V)^*, L^2(I,V)}=0 \quad \forall v \in Q^0(I,V)
\end{align*}

\end{prop}


We have also seen that with smoothness assumption on data (\cref{ass:diri} and \cref{ass:neu}) we obtain that the solutions of \cref{pb:diri}, \cref{pb:neu} have $Q_0(I,V)$ smoothness. 

We can therefore formulate the two partial differential equations directly on $Q_0(I,V)$ as follows.

\begin{align*}
w \in W_0(I, H^1_c),\bar{u}+v_0 \in W_0(I,H^1_c), v_0 \in W_0(I,H^1_0)\\
w' + A w = (g,\cdot)_{L^2(\Gamma_f)} \text{ in }H^{1*}_c \text{ and for a.e. } t \in (0,T) \\
v_0' + A v_0 = -((\bar{u}',\cdot)_H+A \bar{u}) \text{ in }H^{-1} \text{ and for a.e. } t \in (0,T) 
\end{align*}

with $\bar{u}$ any given $\bar{u}\in H^1(I,H^1_c)$ such that $\tr \bar{u} =f$ on $\Sigma_f$, and with $\bar{u}(0)=0$.

We are working under \ref{ass:diri}, \ref{ass:neu}.

Thanks to \cref{prop:eq_form}, this is equivalent to:

\begin{align*}
w \in W_0(I, H^1_c), \bar{u}+v_0 \in W_0(I,H^1_c), v_0 \in W_0(I,H^1_0) \\
\int_I \langle w' , q\rangle_{H^{1*}_c,H^1_c}+ (\nabla w, \nabla q)_H = \int_I(g,\tr q)_{L^2(\Gamma_f)}, \quad \forall q \in Q^0(I, H^1_c) \\
\int_I \langle v_0',p\rangle_{H^{-1},H^1_0} + (\nabla v_0, \nabla p)_H= -\int_I(\bar{u}',p)_H+(\nabla \bar{u}, \nabla p)_H, \quad \forall p \in Q^0(I, H^1_0) 
\end{align*}

By regularity, see \cref{prop:diri_wp}, \cref{prop:wp_neu}, and thanks to \cref{prop:Q} this implies:
\begin{align*}
w \in Q_0(I, H^1_c), \bar{u}+v_0 \in Q_0(I,H^1_c), v_0 \in Q_0(I,H^1_0) \\
\int_I ( w' , q)_H+ (\nabla w, \nabla q)_H = \int_I(g,\tr q)_{L^2(\Gamma_f)}, \quad \forall q \in Q^0(I, H^1_c) \\
\int_I (v_0',p)_H + (\nabla v_0, \nabla p)_H= -\int_I(\bar{u}',p)_H+(\nabla \bar{u}, \nabla p)_H, \quad \forall p \in Q^0(I, H^1_0) 
\end{align*}

where now the derivatives are in the $H^1(I,H)$ sense. Indeed we have proved in \cref{prop:Q} that $u \in W(I,V)$ with $L^2(I,H)$ derivative, is actually $Q(I,V)$, with weak derivative in the $H^1(I,H)$ sense equal to the $L^2(I,H)$ representative of $u'$ in the $W(I,V)$ sense. There, we also proved the representation of the duality bracket.

Conversely, a solution $w \in Q_0(I, H^1_c), \bar{u}+v_0 \in Q_0(I,H^1_c), v_0 \in Q_0(I,H^1_0) $ to the above problem satisfies $w \in W_0(I, H^1_c), \bar{u}+v_0 \in W_0(I,H^1_c), v_0 \in W_0(I,H^1_0)$, see \cref{prop:Q}, and the proof of \cref{prop:diri_wp} . And by \cref{prop:Q} we can get to:

\begin{align*}
w \in W_0(I, H^1_c), \bar{u}+v_0 \in W_0(I,H^1_c), v_0 \in W_0(I,H^1_0) \\
\int_I \langle w' , q\rangle_{H^{1*}_c,H^1_c}+ (\nabla w, \nabla q)_H = \int_I(g,\tr q)_{L^2(\Gamma_f)}, \quad \forall q \in Q^0(I, H^1_c) \\
\int_I \langle v_0',p\rangle_{H^{-1},H^1_0} + (\nabla v_0, \nabla p)_H= -\int_I(\bar{u}',p)_H+(\nabla \bar{u}, \nabla p)_H, \quad \forall p \in Q^0(I, H^1_0) 
\end{align*}

By \cref{prop:eq_test} we obtain back:

\begin{align*}
w \in W_0(I, H^1_c),\bar{u}+v_0 \in W_0(I,H^1_c), v_0 \in W_0(I,H^1_0)\\
w' + A w = (g,\cdot)_{L^2(\Gamma_f)} \text{ in }H^{1*}_c \text{ and for a.e. } t \in (0,T) \\
v_0' + A v_0 = -((\bar{u}',\cdot)_H+A \bar{u}) \text{ in }H^{-1} \text{ and for a.e. } t \in (0,T) 
\end{align*}

Therefore:

\begin{prop}[Equivalent formulation]
\label{prop:eq_form}

Under \cref{ass:diri}, \cref{ass:neu}, \cref{pb:diri}, \cref{pb:neu} can be equivalently formulated as:

\begin{align*}
w \in Q_0(I, H^1_c), \bar{u}+v_0 \in Q_0(I,H^1_c), v_0 \in Q_0(I,H^1_0) \\
\int_I ( w' , q)_H+ (\nabla w, \nabla q)_H = \int_I(g,\tr q)_{L^2(\Gamma_f)}, \quad \forall q \in Q^0(I, H^1_c) \\
\int_I (v_0',p)_H + (\nabla v_0, \nabla p)_H= -\int_I(\bar{u}',p)_H+(\nabla \bar{u}, \nabla p)_H, \quad \forall p \in Q^0(I, H^1_0) 
\end{align*}

Existence, uniqueness and stability proved already in  \cref{prop:diri_wp}, \cref{prop:wp_neu} carry over to this new formulation.

\end{prop}

\chapter{Domains transformations}

\section{Transforming domains}

\begin{prop}[Measurability of composition]
\label{prop:circ_wd}

Define $\pazocal{M}:=\{\tau: \mR^n \rightarrow \mR^n \text{ Lebesgue measurable } \}/\sim$, the quotient being the almost everywhere equal relation (according to the Lebesgue measure).

Consider also $U$ from $\pazocal M_c :=\{\tau: \mR^n \rightarrow \mR^n \text{ continuous } \}/\sim$, the application "unique continuous representative", and $\pazocal{M}_{BL} :=\{\tau: \mR^n \rightarrow \mR^n \text{ Lipschitz homeomorphism} \}/\sim \subseteq \pazocal{M}_c$.

We can then define $\circ:  \pazocal{M} \times \pazocal{M}_{BL}, \pazocal{M}_c \times \pazocal{M} \rightarrow \pazocal{M}$ by, respectively, $[f]\circ g := [f\circ U(g)], f \circ [g]:= [U(f)\circ g]$. These definitions are well posed.

\end{prop}
\begin{mproof}

\underline{$\pazocal{M}_c$}

$U(f)$ is Borel measurable, so the preimage of a Borel set is Borel measurable, and $g$ is Lebesgue measurable, so his preimage of such Borel set is Lebesgue measurable (\textcolor{red}{see \href{https://math.stackexchange.com/questions/283443/is-composition-of-measurable-functions-measurable}{here} for the different notions of measurability}).

This shows that $U(f)\circ g$ is measurable.

To complete the well posedness, if $h=g$ a.e., then clearly $U(f)\circ g= U(f)\circ  h$ a.e..

\underline{$\pazocal{M}_{BL}$}

Consider $f\circ U(g)$. We need to prove it is measurable and that is only depends on $[f]$.

For the measurability: the preimage of a Borel set, by $f$, is Lebesgue measurable $L$. $U(g)$ has a Lipschitz inverse, which will map this set to a Lebesgue set. Indeed, $L = B \cup N$, with $B$ Borel and $N$ Lebesgue measurable and null (\textcolor{red}{see \href{https://math.stackexchange.com/questions/3420145/lebesgue-measurable-set-union-of-borel-set-and-null-set}{here}}). Image and unions commute, so, $U(g)^{-1}(L) = U(g)^{-1}(B) \cup U(g)^{-1}(N)$. The first set is Lebesgue measurable by measurability of $U(g)$, the second one is null, because Lipschitz maps map null sets into null sets, see 9.54 at page 271 of \cite{leoni}. 

%Suppose $f,g \in [f]$. Then $f=g$ everywhere but on the null set $E$. Because $\psi \in \cT^1$ we know that $U(\psi)$ is a Lipschitz homeomorphism, so that $U(\psi)^{-1}(E)$ has zero measure by the lemma of Vitali. For the same reason, the composition is measurable, see \cite{murat}, remarque 2.2, page II-7.

\end{mproof}


Throughout, $D$ is a bounded Lipschitz domain. We define as in \cite{murat} the following spaces of transformations:

\begin{defn}[Spaces of transformations]
We define:
\begin{itemize}	
	\item $\cV^k=\{\tau \in \pazocal{M}, \tau-\id \in W^{k,\infty}(\mR^n,\mR^n)\}$, $k\geq 1$
%	\item $\circ: \cV^1 \times \cV^1 \rightarrow \cV^1$ and $\circ:W^{1,\infty}(\mR^n,\mR^n) \times \cV^1 \rightarrow W^{1,\infty}(\mR^n,\mR^n)$ by $\tau_1\circ\tau_2 = [U(\tau_1)\circ \tau]$, for any $[\tau]=\tau_2$, $[]$ being an equivalence class according to $\sim$.
	\item $\cT^k=\{\tau \in \cV^k \text{ with an } \eta \in \cV^k, \tau \circ \eta = \eta \circ \tau = \id\}$. Any such $\eta$ is unique, we denote it by $\tau^{-1}$ and we have that $U(\tau)$ is a Lipschitz homeomorphism with $U(\tau^{-1})=U(\tau)^{-1}$
\end{itemize} 
\end{defn}

\begin{obs}[A technicality]

Technically, in the original definition of \cite{murat}, $\tau$ need not to be a continuous function, although this is suggested e.g. in remarque 2.1 at page II-4. 

Going to equivalence classes of $\tau$ makes the identification with continuous functions more precise, as we now show.

\underline{One implication}

Let $\tau: \mR^n\rightarrow\mR^n$ with $[\tau-\id] \in W^{k,\infty}$. Then $\tau$ is equal a.e. to a (Lebesgue) measurable function, hence also (Lebesgue) measurable, and thus $[\tau] \in \cV^k$ as we have defined it (\textcolor{red}{this is proved \href{https://heil.math.gatech.edu/6337/spring11/section3.4.pdf}{here}; note that $\{g\neq f\}$ is measurable as the Lebesgue measure is complete}).

Now, suppose $\tau$ is a bijection, and $[\tau^{-1}-\id] \in W^{k,\infty}$ too. Then $\tau = \id + g = G, \tau^{-1} = \id + h = H$ almost everywhere. Here, $G,H$ are at least Lipschitz.
But then $\tau \circ H = \id $ a.e., and since $H$ is Lipschitz, we can conclude also $G\circ H = \id$ a.e., so, everywhere. With a symmetric reasoning, we are lead to $G=H^{-1}$, so that $G$ is bi-Lipschitz.

Thus, $[\tau]\circ [\tau^{-1}]:=[U(\tau)\circ U(\tau^{-1})] = [G\circ G^{-1}] = \id$ and an analogous reasoning leads to $[\tau] \in \cT^k$ as we have defined it.

\underline{The other implication}

It is immediate for $\cV^k$ and for $\cT^k$, in the equivalence class of $\tau \in \cT^k$ there is a unique $U(\tau)$ at least bi-Lipschitz, hence invertible, with $[U(\tau)]=\tau$.

This shows that:

\begin{enumerate}
\item $\{\tau: \mR^n\rightarrow\mR^n$ with $[\tau-\id] \in W^{k,\infty}\} / \sim = \cV^k$
\item $\{\tau: \mR^n\rightarrow\mR^n$ bijection with $[\tau^{\pm 1}-\id] \in W^{k,\infty}\}/\sim = \cT^k$
\end{enumerate}

\end{obs}

We need to check the well-posedness of $\circ$.

\begin{prop}
\label{prop:circ_wd_V}
$\circ: \cV^1 \times \cV^1 \rightarrow \cV^1$ and $\circ:W^{1,\infty}(\mR^n,\mR^n) \times \cV^1 \rightarrow W^{1,\infty}(\mR^n,\mR^n)$ are well defined.
\end{prop}
\begin{mproof}

We start by $\circ:W^{1,\infty}(\mR^n,\mR^n) \times \cV^1 \rightarrow \cV^1$. We have $\te\circ \tau =[ U(\te)\circ U(\tau)]$ for instance (see \cref{prop:circ_wd}); the latter is a bounded Lipschitz map, so it remains in $W^{1,\infty}$.

For the second claim, just write $\eta \circ \tau -\id = (\eta - \id)\circ \tau + \tau -\id$ and use the first part. 

\end{mproof}


\begin{prop}[Chain rule for $k=1$]
\label{prop:chain}
Let $f \in W^{1,\infty}(\mR^n,\mR^n)$ or $\cV^1$, together with $\psi \in \cT^1$.  Then:

\begin{itemize}

\item $f \circ \psi$ has essentially bounded weak derivatives, and $D(f \circ \psi) = Df \circ \psi D\psi$. The equality holds a.e. also for the classical derivatives.

\item  $D(\psi^{-1}) = (D\psi)^{-1} \circ \psi^{-1}$, where $(D(\psi^{-1}))^{-1}:=[(DU(\psi^{-1}))^{-1}]$, the representative being a.e. invertible. The equality holds a.e. also for the classical derivatives.

\item $|\det(D\psi)|$ is an essentially bounded measurable function with $|\det(D\psi)|\geq \delta>0$ a.e.. 
\end{itemize} 

\end{prop}
\begin{mproof}

\underline{Weak derivatives}

We notice that $f \circ \phi$ has a unique Lipschitz representative, that is $U(f)\circ U(\phi)$. The desired formula follows as in \cite{murat}, lemme 2.1 at page II-6, for the classical derivatives, because Lipschitz function are almost everywhere differentiable by the Rademacher theorem (\tred{see \href{https://abel.math.harvard.edu/archive/212b_spring_05/handouts/Rademacher.pdf}{here}}). The chain rule holds for functions differentiable only at one point. 

% (or from \cite{ziemer}, page 53, theorem 2.2.2 in the case of left composition by $ W^{1,\infty}(\mR^n,\mR^n)$ vector fields) (	\tred{this needs the Rademacher theorem, see \href{https://abel.math.harvard.edu/archive/212b_spring_05/handouts/Rademacher.pdf}{here}. The set of differentiability is measurable, see \cite{leoni}, 9.17}).

Now, to identify the weak derivatives:

\begin{itemize}
	\item $U(f)$ is Lipschitz, so that $DU(f)$, the classical derivative, is also the weak derivative $Df$ (note that $f$ need not to be essentially bounded to state this). The latter is a measurable function, as a.e. limit of difference quotients.
%	\item $DU(f)\circ U(\psi)$,  as pointed out in \cite{murat}, remarque 2.2, page II-7, is measurable. It is also essentially bounded.
	\item $DU(f)\circ U(\psi)$, is measurable, see \cref{prop:circ_wd}. It is also essentially bounded.
	\item By \cref{prop:circ_wd} we observe that $DU(f)\circ U(\psi)$ represents $Df \circ \psi$
	\item $D\psi = [DU(\psi)]$ as seen above
	\item the product of equivalence classes is always defined as the product of their representatives
\end{itemize}

Therefore $ Df \circ \psi D\psi = [DU(f)\circ U(\psi) DU(\psi)]$.

And now, because $f \circ \phi$ is Lipschitz, it has weak derivatives, $D(f \circ \phi)$, equal to the classical derivatives $DU(f\circ \phi) = D (U(f)\circ U(\psi)) = DU(f)\circ U(\psi) DU(\psi)$, where the last equality holds a.e., as mentioned at the beginning of the proof.

This let us conclude the first claim.

\underline{Inverse Jacobian}

For the second one, put $f = \psi^{-1}$. Then, for the classical derivatives, $I = DU(\psi)\circ U(\psi)^{-1} DU(\psi^{-1})$ a.e., so that both $DU(\psi)\circ U(\psi)^{-1}, DU((\psi)^{-1})$ are invertible as matrices, a.e.. 

\underline{Determinant}

We have defined $|\det(D\psi)|:=[|\det DU(\psi)|]$, see \cref{prop:circ_wd}. The claim follows as in lemme 4.2, pag. IV-7 of \cite{murat}, and because $\det$ is a polynomial of essentially bounded functions.

\end{mproof}

We go on to define the space of admissible transformations.

\begin{defn}[Admissible transformations]
\label{def:adm}
We define $\Theta:=\{\theta \in W^{1,\infty}(\mR^n,\mR^n) \text{ with } \theta=0 \text{ on } \mR^n \setminus D\}$, a Banach subspace of $ W^{1,\infty}(\mR^n,\mR^n)$.


We also define $\cT:=\{\tau \in \cT^1, \tau^{\pm 1}|_{\mR^n\setminus D}=\id\}$. 

\end{defn}

\begin{prop}[Some group properties of $\cT$]
\label{prop:group}
Let $\eta, \tau \in \cT, \te \in \Te$. Then:

\begin{itemize}
	\item $\eta \circ \tau \in \cT$
	\item $\te \circ \tau \in \Te$
	\item $\id$ is the neutral element
	\item $\eta^{-1} \in \cT$
\end{itemize}

\end{prop}
\begin{mproof}

%\underline{Stability under composition (regularity)}
%
%We start by showing that $\tau \circ \eta |_D \in W^{2,\infty}$.
%
%%By $\tau \in \cV^1$, we sure have that $\tau \circ \eta \in \cV^1$ by \cref{prop:circ_wd_V}. So, 
%
%$U(\tau)\circ U(\eta)$, is a bounded Lipschitz function. Morevoer $ D(U(\tau)\circ U(\eta)) = D(U(\tau))\circ U(\eta) DU(\eta) $ everywhere in $D$, because $\eta, \tau$ happen to be in $W^{2,\infty}(D)$, so that $U(\eta), U(\tau) \in C^{1,1}(D)$, see \cref{prop:lip}.
%
%This is a product of a bounded Lipschitz function and a bounded Lipschitz function, so it is also bounded Lipschitz. Therefore, $D(U(\tau)\circ U(\eta))\in C^{0,1}(D)$ and thus $U(\tau)\circ U(\eta) \in C^{1,1}_B(D)$, so that by $\ref{prop:lipk}$, $[U(\tau)\circ U(\eta)|_D] = \tau \circ \eta|_D \in W^{2,\infty}(D)$.
%
%This same proof shows that for $\te \in \Te$, $\tau \in \cT$, $\te \circ \tau \in \Te$, because $\te \circ \tau$ was already a $W^{1,\infty}$ function by \cref{prop:circ_wd}, and because $\tau$ fixes $\mR^n \setminus D$.

\underline{Stability under inversion}

It is trivial, because the definition of $\cT$ is symmetric with respect to inversion.

\underline{Stability under composition ($\cT^1$)}

$\eta \circ \tau$ is surely in $\cV^1$ by \cref{prop:circ_wd}. Now, by the above point, $\tau^{-1} \circ \eta^{-1}$ is in $\cV^1$ too, and the composition yields: $(\eta \circ \tau)\circ (\tau^{-1} \circ \eta^{-1}) = [U(\eta)\circ U(\tau)]\circ [(U\tau)^{-1} \circ (U\eta)^{-1}] = \id$.

\end{mproof}

\begin{prop}[Small perturbations of $\cT$]
\label{prop:ptb_id}
Let $\te \in \Te$ with small enough $\norm{\te}_{W^{1,\infty}(\mR^n;\mR^n)}$. Then, $\id+\te\in \cT$.

Let $\delta \te \in \Te$ with small enough $\norm{\delta\te}_{W^{1,\infty}(\mR^n;\mR^n)}$, and $\tau \in \cT$. Then, $\tau + \delta \te\in \cT$.

\end{prop}
\begin{mproof}

\underline{Perturbation of identity}

We only need to check the properties of the inverse map.

$U(\tau)^{-1}$ exists and is Lipschitz, see the proof of lemme 2.4 of \cite{murat}, page II-16. We can therefore define $\tau^{-1}$ and we obtain that it is $\cV^1$. So, $\tau \in \cT^1$. The fact that $\tau = \id$ outside of $D$ automatically implies $\tau^{-1}=\id$ outside of $D$, which can be seen more precisely by going to the smooth representatives of $\tau, \tau^{-1}$.

%Moreover, $D\tau$ is a function in $W^{1,\infty}(D;\mR^{n\times n})$. Because of the Banach algebra properties listed in proposition 2.1 of \cite{murat}, page II-5 (which apply to $D$, just by extending $W^{1,\infty}(D)$ functions to $0$), the theorem of Neumann series holds (for $D\te$!), and by the assumptions on the norm of $\te$ we observe that $D\te$ has a small norm $W^{1,\infty}(D;\mR^{n\times n})$: by $D\tau   = I + D\te$ , we conclude that there is $ S \in W^{1,\infty}(D;\mR^{n\times n})$ with $SD\tau = D\tau S = I$, in the sense of matrix multiplication.
%
%This shows $(D\tau)^{-1} \in W^{1,\infty}(D;\mR^{n\times n})$.
%
%Because $\tau \in \cT^1$, we have from \cref{prop:chain}, that $D\tau^{-1} = (D\tau)^{-1}\circ \tau^{-1}$, so that $D\tau^{-1} $ is the composition of a $W^{1,\infty}(D;\mR^{n\times n})$ by a $\cT$ function. This implies that $\tau^{-1}$ is $W^{2,\infty}(D)$, because $\tau^{-1}$ maps $D$ into $D$.

\underline{Perturbation, not of identity}

We solve the equation $\tau + \delta \te =\eta \circ \tau$, i.e., we define $\eta:=\id + \delta \te \circ \tau^{-1}$. Because $\tau^{-1} \in \cT$ and $\delta \te \in \Te$ we observe that $\delta \te \circ \tau^{-1} \in \Te $, thanks to \cref{prop:group}.

We only need to prove that $\delta \te \circ \tau^{-1}$ is small, and then use the first part. 

But by \cref{prop:chain} this follows immediately. One can alternatively apply point i) of lemme 2.2, \cite{murat}.

\end{mproof}

\begin{thm}{Small perturbations of identity, Lipschitz property}
\label{thm:ptb_id_lip}
Let $U\cc D$ be Lipschitz bounded. There exists $0<C(U)<1$ such that, for $\tau \in W^{1,\infty}(\mR^n;\mR^n)$ and $\norm{\tau - \id}_{W^{1,\infty}(\mR^n;\mR^n)}\leq C(U)$, then $T(U)$ is also bounded Lipschitz, where $T$ is $U(\tau)$, the unique Lipschitz continuous representative of $\tau$ (see \cref{prop:lip}).

This result can be applied to, e.g., $\tau \in \cT$ which is a small perturbation of identity in the $W^{1,\infty}$ topology.
\end{thm}

\begin{mproof}
It is done in \cite{bello}, lemma 3, page 629.
\end{mproof}

%\begin{prop}[Gateaux differentiability]
%Consider $J:\cT \rightarrow E$ for $E$ some Banach space.
%
%Let $\tau \in \cT$. Then:
%
%$$ \forall \delta\te \in \Te \text{ exists } \lim_{t\rightarrow 0}\frac{J(\tau + t \delta\te)-J(\tau)}{t} \iff  \forall \delta\te \in \Te  \text{ exists }\lim_{t\rightarrow 0}\frac{J((\id +t \delta \te)\circ \tau)-J(\tau)}{t}$$
%
%In case of existence, we have:
%
%$$ \lim_{t\rightarrow 0}\frac{J((\id +t \delta \te)\circ \tau)-J(\tau)}{t} = \lim_{t\rightarrow 0}\frac{J(\tau + t \delta \te\circ \tau)-J(\tau)}{t}$$
%
%Note that $\id + t \delta \te \in \cT$ for small $t$.
%
%\end{prop}
%
%\begin{mproof}
%
%It suffices to show that $\Te \circ \tau = \Te$.
%
%In fact, $\delta \te \circ \tau \in \Te$ for $\te \in \Te$, as we verified in \cref{prop:group}. Using the same result, $\tau^{-1} \in \cT$, $(\delta \te \circ \tau ) \circ \tau^{-1}  \in \Te$ and $(\delta \te \circ \tau ) \circ \tau^{-1} = [U(\delta \te)\circ U(\tau)]\circ[U(\tau)^{-1}] = [(U(\delta \te)\circ U(\tau))\circ(U(\tau)^{-1})]=\delta \te$.
%
%\end{mproof}
%

\section{Transforming Sobolev spaces}

\begin{thm}[Change of variables]
\label{thm:change}

Let $U$ be open and $T = U(\tau)$ for $\tau \in \cT^1$, and let $p \in [ 1,\infty]$. Then:

\begin{enumerate}
	\item $f \in L^p(T(U)) \iff f\circ T \in L^p(U)$ and there holds, for $f \in L^p(T(U))$:
	$$ \norm{f}_{L^p(T(Q))}\leq \left ( \norm{\det DT}_{L^\infty(\mR^n)}\right)^{1/p} \norm{f\circ T}_{L^p(Q)}$$
	\item $f \in W^{1,p}(T(U)) \iff f\circ T \in W^{1,p}(U)$ and there holds, for $f\in W^{1,p}(T(U))$:
	$$Df \circ T = (Df)^{-t}D(f\circ T)$$
	$$ \norm{Df}_{L^p(T(Q);\mR^n)}\leq \left ( \norm{\det DT}_{L^\infty(\mR^n)}\right)^{1/p} \norm{(DT)^{-1}}_{L^\infty(\mR^n;\mR^{n\times n})}\norm{D(f\circ T)}_{L^p(Q;\mR^n)}$$
	\item \tred{add the rest of this proposition}
	\item if $p \in (1, \infty)$, $f \in W^{1,p}_0(T(U)) \iff f\circ T \in W^{1,p}_0(U)$
	\item therefore, composition by $T$ is a linear isomorphism between $W^{k,p}(T(U))\rightarrow W^{k,p}(U)$ for $k=0,1$, and between $W^{1,p}_0(T(U))\rightarrow W^{1,p}_0(U)$ for $k=0,1$, $p \in (1, \infty)$
	\item for $D$ a bounded Lipschitz domain and $\cT, \Te$ defined before, we get, for $f \in H^1(D)$, that $\tr f = \tr(f\circ T)$
	\item if moreover, $\Omega, T(\Omega) \cc D$ are also bounded Lipschitz domains, letting $U:=D\setminus \Omega$, another bounded Lipschitz domain, for $f \in H^1(T(U))$ and $\tr_{T(U)} f =0 $ on $ \partial T(\Omega) $, then $\tr_{U} f\circ T=0$ on $\partial \Omega$ and $\tr_{T(U)} f = \tr_{U} f\circ T$ on $\partial D$
	\item so, $\circ T$ is a linear isomorphism of $H^1_c(U)$ and $H^1_c(T(U))$ ($H^1_c$ is defined in \cref{subs:inh_diri})
\end{enumerate}

\end{thm}

\begin{mproof}

We need to prove only the last points, for the other are proved in \cite{murat}, pages IV.4, IV.5, IV.6.

\underline{Static strace}

To do so, let $f_n \in C(\overline{D})\cap H^1(D)$ converging in $H^1(D)$ to $f$. By point 4, we have $f_n \circ T\rightarrow f\circ T$ in $H^1(D)$ (rememeber, $T(D)=D$ by invertibility of $T$ and the fact that $T(x)=x$ outside of $D$). Therefore we have:

$$\tr f \leftarrow_{L^2(\partial D)}\tr(f_n) = f_n|_{\partial D} = (f_n\circ T)|_{\partial D}= \tr(f_n \circ T)\rightarrow_{L^2(\partial D)} \tr(f \circ T)$$

\underline{Zero moving trace}

First of all, as $T$ is a homeomorphism of $\mR^{n}$, $TU=D\setminus T(\Omega)$, $T\partial U = \partial D \sqcup \partial \Omega$, $T\partial \Omega = \partial T \Omega$.

Now, an application of \cref{thm:ibp} yields that the extension to $0$ in $T\Omega$ of $f$, call it $\bar{f}$, is $H^1(D)$, with $\partial_i \bar{f}=\partial_i f$ in $TU$, $0$ in $T(\Omega)$.

We claim that $\tr_D \bar{f} = \tr_{T(U)} f|_{\partial D}$. In fact, approximate $\bar{f}$ by restrictions to $D$ of $C^\infty_c(\mR^n)$ functions $f_n$ (see theorem 3.18 of \cite{adams}, page 54), which also approximate $f$ on ${T(U)}$, by the observation that $\norm{f_n|_{T(U)}}_{H^1({T(U)})}\leq \norm{f_n|_D}_{H^1(D)}$. Then:

$$\tr_{T(U)}( f_n|_{T(U)})|_{\partial D} = (f_n|_{T(U)})|_{\partial {T(U)}}|_{\partial D} = f_n|_{\partial D} = \tr_D ( f_n|_D ) $$

Now, by what we observed before, $\tr_{T(U)}( f_n|_{T(U)})\rightarrow \tr_{T(U)}(f) $ in $L^2(\partial {T(U)})$, so that $\tr_{T(U)}( f_n|_{T(U)})|_{\partial D} \rightarrow \tr_{T(U)}( f)|_{\partial D}$. On the other hand $\tr_D ( f_n|_D ) \rightarrow \tr_D \bar{f}$, which yields the claim.

Using this: $ \tr_{T(U)}( f)|_{\partial D} = \tr_D \bar{f} = \ind{point 5} = \tr_D(\bar{f}\circ T) = \tr_D(\overline{f \circ T}) = \tr_U (f\circ T)|_{\partial D}$, where we used that $\bar{f}\circ T$ is zero in $T^{-1}T\Omega = \Omega$ (because again $T$ maps null sets into null sets), so it is the zero extension $\overline{f \circ T}$ of $f\circ T$, and applied the same reasoning as above to conclude $\tr_D(\overline{f \circ T}) = \tr_U (f\circ T)|_{\partial D}$. Both $\bar{f}\circ T$ and $f\circ T$ are $H^1$ functions by point 2.

We can now also say that $\tr_U f\circ T=0 $ on $\partial \Omega$.

%In fact, let $\phi \in C^\infty_c(D)$. Then $ \int_D \bar{f}\circ T \phi_{,i} = \int_U  f\circ T \phi_{,i} = \ind{using \cref{thm:ibp}} = \int_{\partial \Omega}\tr_U(f\circ T) \phi \nu_i d\sigma -\int_U \partial_i(f\circ T)\phi =  \int_{\partial \Omega}\tr_U(f\circ T) \phi \nu_i d\sigma -\int_D \partial_i(\bar{f}\circ T)\phi  =  \int_{\partial \Omega}\tr_U(f\circ T) \phi \nu_i d\sigma +\int_D \bar{f}\circ T \phi_{,i}  $.
%
%It follows that $\int_{\partial \Omega}\tr_U(f\circ T) \phi \nu_i d\sigma = 0$ for all $\phi \in C^\infty_c(D)$. Let now $\phi_n \in C^\infty_c(\mR^n)$ approximate $f\circ T$ in $H^1(U)$ as seen above. Choose a cut-off function $\eta$ that is $1$ in a small neighbourhood of $ \Omega$ and $0$ in the complement of $D$ and a small neighbourhood of $\partial_D$. Then $\eta \phi_n$ is in $C^\infty_c(D)$ and:

$$(\eta \phi_n)|_{\partial \Omega} = \tr_U( \phi_n|_U)_{\partial \Omega} \rightarrow \tr_U(f \circ T)_{\partial \Omega}$$
%
%the limit being in $L^2(\partial \Omega)$. Because  $\eta \phi_n$ is in $C^\infty_c(D)$ we can choose it $\phi$ in $\int_{\partial \Omega}\tr_U(f\circ T) \phi \nu_i d\sigma = 0$, yielding:
%
%$$\int_{\partial \Omega}\tr_U(f\circ T) \phi_n \eta \nu_i d\sigma = 0$$ 
%
%This implies, thanks to the boundedness and measurability of $\nu_i$, that:
%
%$$\int_{\partial \Omega}\tr_U(f\circ T)^2\nu_i d\sigma = 0$$

\underline{Multiplication by a $W^{1,\infty}$ function}

We claim that, for $\psi \in W^{1,\infty}(\mR^n;\mR)$ and $f \in H^1(U)$, then $f\psi$ has the same trace as $f$ as long as $\psi = 1$ in a neighbourhood of $\partial U$.

Note that $f\psi \in H^1(U)$ still. Now: approximate $f$ by restriction of test functions $f_n$. Then $f_n \psi$ is $C(\overline{U})\cap H^1(U)$ (thanks also to \cref{prop:lipk}), so that $\tr_U(f_n\psi) = \tr_U(f_n)$. Because $f_n \psi \rightarrow f \psi$ is $H^1(U)$ the claim is valid.

This last convergence follows from $\norm{(f_n-f)\psi}_{L^2}\leq \norm{(f_n-f)}_{L^2}\norm{\psi}_{L^\infty}$, the chain rule $\partial_i(f_n\psi)=\partial_i f_n \psi + \partial_i \psi f_n$ (see \textcolor{red}{corollary 4.1.18 \href{https://www.math.stonybrook.edu/~joa/PUBLICATIONS/SOBOLEV.pdf}{here}}) and again $\norm{\partial_i(f_n-f)\psi}_{L^2}\leq \norm{\partial_i(f_n-f)}_{L^2}\norm{\psi}_{L^\infty}$, $\norm{(f_n-f)\partial_i\psi}_{L^2}\leq \norm{(f_n-f)}_{L^2}\norm{\partial_i\psi}_{L^\infty}$.

\underline{Reducing to a function of $0$ trace}

Let $\eta$ be a smooth cut-off function which is $1$ close to $\partial D$ and $0$ close to $\partial T\Omega$, $\beta=0$ close to $\partial D$ and $1$ close to $\partial T\Omega$. This can be accomplished by e.g. building a suitable partition of unity of the compact sets $\partial \Omega$ and $\partial D$. (\textcolor{red}{can I do this? Yes, see bachelor's thesis, take $K =\partial \Omega$ etc. Also, be careful with all of these equalities...}).

%Then $\beta + \eta = 1$ on $\partial T(U)$ and thus $\tr_{T(U)} f = \tr_{T(U)} (f(\eta+\beta)) = \tr_{T(U)} f\eta+\tr_{T(U)} f\beta$.

$f\beta$ has zero trace, as it can be verified by approximating $f$ by smooth functions again: 

$$\tr_{T(U)} f\beta \leftarrow_{L^2(\partial T(U))} \tr_{T(U)} f_n \beta $$

where the latter quantity is $\tr_{T(U)} f_n $ on $\partial T(U)$ and $0$ on $\partial D$. By restricting the convergence  to first $\partial D$ and then to $\partial T(U)$, and using almost everywhere convergent subsequences, we conclude that $\tr_{T(U)} f \beta = \tr_{T(U)} f $ on $\partial T(U)$ and  $\tr_{T(U)} f \beta = 0 $ on $\partial D$, i.e.  $f\beta$ has zero trace.

%The same argument yields that $\tr_{T(U)} f \eta = \tr_{T(U)} f $ on $\partial D$ and  $\tr_{T(U)} f \eta = 0 $ on $\partial T(U)$.
%
%Hence $\tr_{T(U)} f |_{\partial T(U)}= \tr_{T(U)} f\beta |_{\partial T(U)}$, with $f\beta$ of zero trace on $T(U)$. 

\underline{Domain transformation}

But zero trace functions in $H^1(T(U))$, since $T(U)$ is assumed to be bounded Lipschitz, are exactly the functions $H^1_0(T(U))$ (theorem 18.7 at page 595 of \cite{leoni}).

By then point 4, $(f\beta)\circ T \in H^1_0(U)$.

Because $T$ is bi-Lipschitz, we can write $(f\beta)\circ T = f \circ T \beta \circ T$ almost everywhere.

We have that $\beta \circ T +  \eta \circ T$ is $W^{1,\infty}$ and $1$ near $\partial U$. 

So, $\tr_U f\circ T = \tr_U f\circ T(\beta \circ T +  \eta \circ T) = \tr_U (f\circ T\beta \circ T) + \tr_U (f\circ T  \eta \circ T)$.

Approximate $f\circ T$ by $g_n$ smooth as seen above.  Then, $\tr_U (g_n \eta \circ T)$ is $0$ on $\partial \Omega$ and $\tr_U g_n$ on $\partial D$. By selecting an almost everywhere convergent subsequence, we conclude $\tr_U (f\circ T \eta \circ T) = 0$ on $\partial \Omega$. 

Hence  $\tr_U f\circ T|_{\partial \Omega} = tr_U f\circ T(\beta \circ T)|_{\partial \Omega} = 0$.

\end{mproof}

%We need to define the pullback of functionals. 
%
%\begin{ass}[Defining pullbacks]
%\label{ass:pull}
%
%Throughout, $\Omega\cc D$ is another bounded Lipschitz domain, and we assume $\tau \in \cT$ preserves this property, i.e., that $T (\Omega) $ is also bounded Lipschitz, $T=U(\tau)$.
%
%We denote by $V$ either $H^1_0(U)$ or $H^1_c(U)$ (see \cref{subs:inh_diri} for the definition of this space, and by $V_T$ either  $H^1_0(T(U))$ or $H^1_c(T(U))$.
%\end{ass}
%
%\begin{prop}[Pullback]
%\label{prop:pull}
%
%Under \cref{ass:pull}, let $F \in V_T^*$. We define $F\circ T \in V^*$ by:
%
%$$\langle F, v\rangle_{V^*_T,V_T} = \langle F\circ T , |\det(DT)| v\circ T\rangle_{V^*,V}$$
%
%The definition is well posed.
%
%Moreover:
%\begin{itemize}
%\item the maps $V\rightarrow V$, $w\mapsto w|\det(DT)|$, and $S_T: V_T\rightarrow V$, $v \mapsto |\det(DT)| v\circ T$ are linear isomorphisms
%\item taking pullbacks, seen as a map $P_T: V_T^*\rightarrow V^*$, is linear and bounded
%\item we have $P_T = S_T^{-*}$, so that taking pullbacks is also a linear isomorphism
%
%\end{itemize}
%
%
%\end{prop}
%\begin{mproof}
%
%\underline{Existence by density}
%
%Let $H_T \ni f_n \rightarrow F$ in $V_T^*$, possible by density of $H_T=L^2(T(U))$ in $V_T^*$.
%
%By \cref{thm:change} we observe that $g_n:=f_n\circ T \in H$ is Cauchy in $V^*$. In fact:
%
%\begin{align*}
%\langle g_n-g_m ,  w \rangle_{V^*,V}  = \int_U (g_n-g_m)w =\ind{change of variables} = \\
%\int_{T(U)} (f_n-f_m)w\circ T^{-1} |\det(D(T^{-1}))| = \\
%\langle f_n-f_m ,  w\circ T^{-1} |\det(D(T^{-1}))| \rangle_{V^*_T,V_T} 
%\end{align*}
%
%We have to prove that $w\circ T^{-1} |\det(D(T^{-1}))| \in V_T$. We start by showing that $|\det(D(T^{-1}))| \in W^{1,\infty}(D)$. Because $\tau \in \cT$, then $\det(D(T^{-1}))\in W^{1,\infty}(D)$, as a product of $W^{1,\infty}(D)$ functions, because we have $\tau^{-1}|_D \in W^{2,\infty}(D)$ (i.e. a product of bounded Lipschitz functions, hence, still bounded Lipschitz, see \cref{prop:lip}). And now, the absolute value of a $W^{1,\infty}(D)$ function, is still $W^{1,\infty}(D)$, because the absolute value preserve the Lipschitz continuity, and the boundedness.
%
%Also, $W^{1,\infty}\subseteq H^1$, so that $w\circ T^{-1} |\det(D(T^{-1}))| \in H^1(T(U))$. For the traces, we know by \cref{thm:change} that $w\circ T^{-1} \in V_T$. Suppose e.g. that $\tr(w\circ T^{-1})=0$ on $\partial T(U)$. Approximate $w\circ T^{-1}$ by $v_n$ smooth to conclude that  $\tr (w\circ T^{-1} |\det(D(T^{-1}))|) = \tr(w\circ T^{-1})\cdot  |\det(D(T^{-1}))|$. Here we have used that a $C^{0,1}_B$ function can be extended by continuity to the boundary, in view of its uniform continuity. This shows that $w\circ T^{-1} |\det(D(T^{-1}))| \in V_T$.
%
%Thus:
%
%\begin{align*}
%\langle g_n-g_m ,  w \rangle_{V^*,V} \leq \\
%\norm{f_n-f_m }_{V^*_T}\norm{  w\circ T^{-1} |\det(D(T^{-1}))| }_{V_T} 
%\end{align*}
%
%Call $d:=|\det(D(T^{-1}))| \in W^{1,\infty}(D)$. By the product rule we get $\norm{D( w\circ T^{-1} d)}_{L^2(T(U))}^2\leq \sum_i \left( \norm{d}_{L^\infty(D)}^2 +\norm{\partial_i d}_{L^\infty(D)}^2\right )^2 \norm{D(w\circ T^{-1})}_{V_T}^2$, which means that:
%
%\begin{align*}
%\langle g_n-g_m ,  w \rangle_{V^*,V} \leq \\
%C(T)\norm{f_n-f_m }_{V^*_T}\norm{w\circ T^{-1} }_{V_T} \leq\\
%C(T)\norm{f_n-f_m }_{V^*_T}\norm{w}_{V}
%\end{align*}
%
%where we used \cref{thm:change} in the last passage. Hence also $g_n$ is Cauchy in the Banach space $V^*$ and we can thus conclude the existence of $G \in V^*$ such that $g_n\rightarrow G$ in $V^*$.
%
%And now:
%
%\begin{align*}
%\langle f,  v \rangle_{V^*_T,V_T} = \lim \langle f_n,  v \rangle_{V^*_T,V_T} = \\
%\lim( f_n,v)_{H_T} =\ind{change of variables} =\\
%\lim( g_n,v\circ T |\det(DT)|)_{H} = \\
%\lim\langle g_n, v\circ T |\det(DT)| \rangle_{V^*,V} = \\
%\langle G, v\circ T |\det(DT)| \rangle_{V^*,V}
%\end{align*}
%
%where we used the fact that $v\circ T |\det(DT)| \in V_T$, in complete analogy with the above reasonings.
%
%\underline{Well posedness}
%
%So, $\langle F,  v \rangle_{V^*_T,V_T} = \langle G, v\circ T |\det(DT)| \rangle_{V^*,V}$.
%
%$|\det(DT)|$ is represented by $d \in C^{0,1}_B(D)$. In \cref{prop:chain} we saw $C\geq d\geq \delta >0$ almost everywhere in $\mR^n$. The inequality holds everywhere on $D$, by continuity of $d$. Then $1/d$ is bounded and Lipschitz, as we can verify by:
%
%$$\left | \frac{1}{d(x)}-\frac{1}{d(y)}\right |=\left| \frac{d(x)-d(y)}{d(x)d(y)}\right |\leq \delta^{-2}C|x-y|$$
%
%Therefore, multiplication by $d$ is a bijection of $V$ (a linear isomorphism actually, we already proved the boundedness, and the boundedness of the inverse follows from the bounded inverse theorem).
%
%Assume there exists another $G_2$ satisfying $\langle F,  v \rangle_{V^*_T,V_T} = \langle G_2, v\circ T |\det(DT)| \rangle_{V^*,V}$.
%
%Then  $\langle G_2-G, v\circ T |\det(DT)| \rangle_{V^*,V} = 0$ for all $v \in V_T$. Because $\circ T$ is another bijection, we can conclude  $\langle G_2-G, w \rangle_{V^*,V} = 0$ for all $w \in V$, which yields the uniqueness.
%
%\underline{Boundedness}
%
%Also, $S_T: V_T \rightarrow V$ defined by $S_Tv = |\det(DT)|v\circ T$ is a linear isomorphism, as a composition of linear isomorphism (see also \cref{prop:chain}). 
%
%Therefore:
%
%$$\langle F, v\rangle_{V^*_T,V_T} = \langle F\circ T , S_Tv\rangle_{V^*,V} =  \langle S_T^*(F\circ T), v\rangle_{V^*,V}$$
%
%This shows that $S_T^{-*}F=F\circ T = P_T F$, and this concludes the proof.
%		
%%Now, $w\circ T^{-1} \in V_T$ by \cref{thm:change}. We can also say that $|w \circ T^{-1} |\in V_T$, with $\norm{|w \circ T^{-1}|}_{V_T}\leq \norm{w \circ T^{-1}}_{V_T}$ (\tred{see \href{https://math.stackexchange.com/questions/2578760/if-u-in-h1-omega-then-u-in-h1-omega}{here} for the bound on the gradient, or \href{https://www.math.ucdavis.edu/~hunter/pdes/ch3.pdf}{here, prop. 3.22}}).
%%
%%In fact, absolute value preserves zero traces. To see this, suppose that $\tr f = 0$ on $\partial \Omega$. Then, as in the proof of \cref{thm:change}, $\tr |f| |_{\partial \Omega} = \tr(|f|\beta)|_{\partial \Omega}$, $\beta$ smooth, $1$ close to $\partial \Omega$ and $0$ close to $\partial D$. Approximating $|f_n|$ by smooth functions we see that: $\tr(|f|\beta)$
%
%\end{mproof}

\section{Transforming Bochner spaces}

\begin{prop}[Isomorphism between $Q$ spaces]
\label{prop:change_boch}

$$\circ T : Q(I,V_T)\rightarrow Q(I,V)$$

is a linear isomorphism, and so is:

$$\circ T : Q_0(I,V_T)\rightarrow Q_0(I,V)$$

In particular, $(u\circ T)' = u'\circ T$ (under suitable identifications, see the proof).

\end{prop}

\begin{mproof}

\underline{Existence of derivative}

Consider the following commutative diagram:

\[\begin{tikzcd}
	{V_T} & {H_T} & {H_T^*} & {V^*_T} \\
	V & H & {H^*} & {V^*}
	\arrow["{a_T}", hook, from=1-1, to=1-2]
	\arrow["{r_T}", from=1-2, to=1-3]
	\arrow["{a_T^*}", hook, from=1-3, to=1-4]
	\arrow["r", from=2-2, to=2-3]
	\arrow["a^*", hook, from=2-3, to=2-4]
	\arrow["a", hook, from=2-1, to=2-2]
	\arrow["{\circ T_H=:t_V}"', from=1-1, to=2-1]
	\arrow["{\circ T_H=:t_H}", from=1-2, to=2-2]
\end{tikzcd}\]

We know that $f\in Q(I,V_T)$ satisfies $a_Tf \in H^1(I,H_T)$. Therefore, by \cref{lemma:bochner_Hk_map}, $t_H a_T f =a t_V f \in H^1(I,H)$, and $t_v  \in L^2(I,V)$, so that $t_V : Q(I,V_T)\rightarrow Q(I,V)$ is well defined.

\underline{Boundedness}

We have that $t_V: V_T \rightarrow V$ is linear bounded, so that $t_V: L^2(I,V_T)\rightarrow L^2(I,V)$ is also linear and bounded. Still by \cref{lemma:bochner_Hk_map}, we have $(a t_V f )' = (t_H a_T f)' = t_H(a_T f)'$.

Thus $\norm{t_V f}_{L^2(I, V)}\leq C(T) \norm{f}_{L^2(I, V_T)}$, together with $\norm{(a t_V f )'}_{L^2(I, H)}\leq C(T) \norm{(a_T f)'}_{L^2(I, H_T)}$.

By noting that $(t_V)^{-1} = (\circ T)^{-1} $ is bijective, and by the bounded inverse theorem:

$$\circ T : Q(I,V_T)\rightarrow Q(I,V)$$

is a linear isomorphism.

\underline{Zero initial value}

Consider $at_V f$. It has a unique $C([0,T],H)$ representative, $U(at_V f)$. Also, $a_T f$ has a unique continuous representative $U(a_T f)$. Now, $at_V f = t_H a_T f$, so that $[U(at_V f)] = [t_H U(a_T f)]$. By continuity, $U(at_V f) = t_H U(a_T f)$ on $[0,T]$ and thus, whenever $U(a_T f)(0)=0$, so is $U(at_V f)$, informally, also $t_Vf (0) =0$. 

So, $t_V(Q_0(I,V_T))\subseteq Q_0(I,V)$.

$(t_V)^{-1} = (\circ T)^{-1}$ and we can conclude that  $t_V^{-1}(Q_0(I,V))\subseteq Q_0(I,V_T)$.

\end{mproof}


\section{Transforming partial differential equations}

We consider again the two parabolic equations of interest, namely, \cref{pb:diri} and \cref{pb:neu}.

We continue from \cref{prop:eq_form}.

\begin{align*}
w \in Q_0(I, H^1_c), v_0 \in Q_0(I,H^1_0) \\
\int_I ( w' , q)_H+ (\nabla w, \nabla q)_H = \int_I(g,\tr q)_{L^2(\Gamma_f)}, \quad \forall q \in Q^0(I, H^1_c) \\
\int_I (v_0',p)_H + (\nabla v_0, \nabla p)_H= -\int_I(\bar{u}',p)_H+(\nabla \bar{u}, \nabla p)_H, \quad \forall p \in Q^0(I, H^1_0) 
\end{align*}

We are working under the assumption:

\begin{ass}
\label{ass:pull}

We have $T=U(\tau), \tau \in \cT$, $U\cc D$ bounded Lipschitz domains and we also assume that $T(U)$ is bounded Lipschitz.
\end{ass}

Suppose the problem is formulated on $T(U)$. To ease the notation, call $H^1_c(T(U))=V_T, H^1_c(U)=V, H^1_0(U)=W$ and analogously for the other spaces.

We write the problem as:

\begin{align*}
w^T \in Q_0(I, V_T), v_0^T \in Q_0(I,W_T) \\
\int_I  (w^T_t , q)_{H_T}+ (\nabla w^T, \nabla q^T)_{H_T} = \int_I(g,\tr_{T(U)} q^T)_{L^2(\Gamma_f)}, \quad \forall q^T \in Q^0(I, V) \\
\int_I (v^T_{0t},p)_{H_T} + (\nabla v_0^T, \nabla p^T)_{H_T}= -\int_I(\bar{u}',p^T)_{H_T}+(\nabla \bar{u}, \nabla p^T)_{H_T}, \quad \forall p^T \in Q^0(I, W_T)
\end{align*}

Applying a change of variables, we get equivalently:

\begin{align*}
w^T \in Q_0(I, V_T), v_0^T \in Q_0(I,W_T) \\
\int_I (w^T_t\circ T, q^T\circ T |\det(DT)|)_H+ (A_T\nabla (w^T\circ T), \nabla( q^T\circ T))_{H} =\\ \int_I(g,\tr_{T(U)} q^T)_{L^2(\Gamma_f)}, \quad \forall q^T \in Q^0(I, V_T) \\
\int_I(v_{0t}^T\circ T,p^T\circ T |\det(DT)|)_H + (A_T \nabla (v_0^T\circ T), \nabla( p^T\circ T))_{H}=\\ -\int_I((\bar{u}')\circ T,p^T\circ T |\det(DT)|)_{H}+(A_T \nabla (\bar{u} \circ T), \nabla (p^T\circ T))_{H}, \quad \forall p^T \in Q^0(I, W_T)
\end{align*}

Here $A_T = |\det(DT)|DT^{-1}(DT)^{-t}$.

Now, we note that:

\begin{itemize}
	\item $\tr_{T(U)} q = \tr_U(q\circ T)$ on $\Sigma_f$ by \cref{thm:change}
	\item $w_t^T\circ T = (w^T\circ T)_t$ by the proof of \cref{prop:change_boch} and analogously for $v_0$
	\item by \cref{prop:change_boch},$\circ T$ is a bijection between $Q^0(I,V_T)$ and $Q^0(I,V)$ and analogously for $W$
	\item $\bar{u} \in H^1(I,V_T)$ and that $\bar{u}'$ denoted the weak derivative in the $H^1(I,V_T)$ sense, so that \ref{prop:bochner_bound} yields $\bar{u}\circ T \in H^1(I,V)$ and $(\bar{u}\circ T )' = \bar{u}'\circ T $
\end{itemize}

We therefore get, equivalently:

\begin{align*}
w^T \in Q_0(I, V_T), v_0^T \in Q_0(I,W_T) \\
\int_I ( (w^T\circ T)_t , q |\det(DT)|)_H+ (A_T\nabla (w^T\circ T), \nabla q)_{H} =\\ \int_I(g,\tr_{U} q)_{L^2(\Gamma_f)}, \quad \forall q \in Q^0(I, V) \\
\int_I ( (v_0^T\circ T)_t,p |\det(DT)|)_H + (A_T \nabla (v_0^T\circ T), \nabla p)_{H}=\\ -\int_I((\bar{u}\circ T)',p|\det(DT)|)_{H}+(A_T \nabla (\bar{u} \circ T), \nabla p)_{H}, \quad \forall p \in Q^0(I, W)
\end{align*}

and by \cref{prop:change_boch}, we also get $w^T\circ T \in Q_0(I,V), v_0^T\circ T \in Q_0(I,W)$.

%\tred{Note, here one could absorb $\det$ into the test functions, complicate the PDE but get rid of the time coefficients. From here one could prove a continuity result for the states and then get Frechet differentiability from Gateaux differentiability}.

On the other hand, consider:

\begin{align*}
w \in Q_0(I, V), v_0 \in Q_0(I,W) \\
\int_I ( w_t , q |\det(DT)|)_H+ (A_T\nabla w, \nabla q)_{H} =\int_I(g,\tr_{U} q)_{L^2(\Gamma_f)}, \quad \forall q \in Q^0(I, V) \\
\int_I ( v_{0t},p |\det(DT)|)_H + (A_T \nabla v_0, \nabla p)_{H}=\\ -\int_I((\bar{u}\circ T)',p|\det(DT)|)_{H}+(A_T \nabla (\bar{u} \circ T), \nabla p)_{H}, \quad \forall p \in Q^0(I, W)
\end{align*}

Then, we note the following:

\begin{itemize}
	\item by \cref{prop:change_boch}, $w\circ T^{-1} \in Q_0(I, V_T), v_0\circ T^{-1} \in Q_0(I,W_T) $, and as seen above, $((w\circ T^{-1})\circ T)_t = (w\circ T^{-1})_t\circ T$ and the same goes for $v_0\circ T^{-1}$
\end{itemize}

Therefore we obtain, equivalently:

\begin{align*}
w\circ T^{-1} \in Q_0(I, V_T), v_0\circ T^{-1} \in Q_0(I,W_T) \\
\int_I ((w\circ T^{-1})_t , q^T)_{H_T}+ (\nabla (w\circ T^{-1}), \nabla q^T)_{H_T} = \\\int_I(g,\tr_{T(U)} q^T)_{L^2(\Gamma_f)}, \quad \forall q^T \in Q^0(I, V_T) \\
\int_I ((v_0\circ T^{-1})_t,p^T)_{H_T} + (\nabla (v_0\circ T^{-1}), \nabla p^T)_{H_T}= \\-\int_I(\bar{u}',p^T)_{H_T}+(\nabla \bar{u}, \nabla p^T)_{H_T}, \quad \forall p^T \in Q^0(I, W_T)
\end{align*}

and $w\circ T^{-1} \in Q_0(I, V_T), v_0\circ T^{-1} \in Q_0(I,W_T)$.

These findings can be summarized as follows.

\begin{thm}[Equivalent formulations with transported domain]
\label{thm:eq_pde}

Let \cref{ass:diri}, \cref{ass:neu}, \cref{ass:pull} hold.

Consider the following problems.

\begin{pb}[Joint parabolic problem, moving domain]
\label{pb:joint_mov}
\begin{align*}
w^T \in Q_0(I, V_T), v_0^T \in Q_0(I,W_T) \\
\int_I  (w^T_t , q^T)_{H_T}+ (\nabla w^T, \nabla q^T)_{H_T} = \int_I(g,\tr_{T(U)} q^T)_{L^2(\Gamma_f)}, \quad \forall q^T \in Q^0(I, V_T) \\
\int_I (v^T_{0t},p^T)_{H_T} + (\nabla v_0^T, \nabla p^T)_{H_T}= -\int_I(\bar{u}',p^T)_{H_T}+(\nabla \bar{u}, \nabla p^T)_{H_T}, \quad \forall p^T \in Q^0(I, W_T)
\end{align*}
\end{pb}

\begin{pb}[Joint parabolic problem, reference domain]
\label{pb:joint_ref}
\begin{align*}
w \in Q_0(I, V), v_0 \in Q_0(I,W) \\
\int_I ( w_t , q |\det(DT)|)_H+ (A_T\nabla w, \nabla q)_{H} =\int_I(g,\tr_{U} q)_{L^2(\Gamma_f)}, \quad \forall q \in Q^0(I, V) \\
\int_I ( v_{0t},p |\det(DT)|)_H + (A_T \nabla v_0, \nabla p)_{H}=\\ -\int_I((\bar{u}\circ T)',p|\det(DT)|)_{H}+(A_T \nabla (\bar{u} \circ T), \nabla p)_{H}, \quad \forall p \in Q^0(I, W)
\end{align*}
\end{pb}

We have the following:

\begin{itemize}
	\item consider $w^T \in Q_0(I, V_T), v_0^T \in Q_0(I,W_T)$. They solve \cref{pb:joint_mov} $\iff$ $w^T\circ T , v_0^T\circ T $ solve \cref{pb:joint_ref}
	\item consider $w \in Q_0(I, V), v_0^T \in Q_0(I,W)$. They solve \cref{pb:joint_ref} $\iff$ $w\circ T^{-1}, v_0\circ T^{-1}$ solve \cref{pb:joint_mov} 

\end{itemize}

Here, $A_T:=  (DT)^{-1}(DT)^{-t}|\det(DT)|$.

\end{thm}


\printbibliography[title={Bibliography}]

\addcontentsline{toc}{chapter}{Bibliography}


\end{document}
